% !TeX root = ../../Book3.tex
% 纲要标题：## 第28章*　平均曲率流、单调性与弱演化
% 纲要位置：几何学计划纲要.md，第 3281 行。
% 正文状态：初稿；尚未正式审核。

% 纲要内容开始
% 
% > 本章目标与证明范围：完成光滑紧初始子流形的短时理论、面积下降及欧氏 Huisken 单调公式。水平集和 Brakke 弱演化有独立定义；奇点后唯一性不由单调量自动得到。
% 
% 纲要内容结束

\OptionalChapter{平均曲率流、单调性与弱演化}{g3:ch:28}
\BookChapterNumber{g3:ch:28}

\begin{chapterabstract}
平均曲率流沿面积的负梯度移动子流形，其参数方程在切向方向退化。本章用初始浸入的法图把方程化为固定向量丛上的强抛物系统，证明闭光滑浸入的短时存在与唯一性，并推导面积耗散。对 Euclid 超曲面，计算第二基本形式的演化，证明避免相交、平均凸性与凸性的保持。随后完整推导任意余维的 Euclid Huisken 单调公式，解释等号、自收缩子及已有光滑放大极限的性质。最后分别定义水平集粘性解与 Brakke 测度演化，以精确的外部存在定理和具体模型辨明集合、重数、质量丢失与肥化。
\end{chapterabstract}

\begin{learninggoals}
	\begin{itemize}
		\item 区分平均曲率法向速度与纯法向参数方程，并从法图推导强抛物系统。
		\item 由第26章的拟线性理论构造短时流，证明固定初始映射后的唯一性。
		\item 推导面积、第二基本形式与平均曲率的演化，识别比较原理的余维及边界条件。
		\item 完成高斯加权面积的微分与分部积分，解释单调公式中的平方项。
		\item 计算球面、圆柱及自收缩模型，分析抛物放大中的空间与时间尺度。
		\item 写出水平集与 Brakke 流的独立定义，判断相应存在性、唯一性和正则性结论的适用对象。
	\end{itemize}
\end{learninggoals}

半径为 \(R\) 的圆球面处处向内弯曲。如果让每一点沿第二基本形式的迹移动，球面保持圆形，而半径满足 \(R'=-m/R\)，所以 \(R(t)^2=R(0)^2-2mt\)。面积不断减小，曲率却不断增大，并在有限时间内出现奇点。这个最简单的模型同时提出三个问题：一般初始子流形是否有这样的运动；哪些几何性质在运动中保留；当光滑描述失效后，应该用什么对象继续描述演化。

\ProseParagraphBreak
第19章的一阶变分把曲率与面积联系起来，第26章建立了求解强抛物方程的工具。两者之间还隔着参数自由：切向移动只是在同一几何形状上重新标记点，方程的主符号因而有核。本章先用初始浸入的法丛表示邻近形状，再把切向速度通过常微分方程消去。这样得到的短时定理适用于一般黎曼环境和任意余维；后续的精确演化、比较及高斯单调公式则在 Euclid 环境中展开。

\ProseParagraphBreak
观察奇点需要在空间中放大，同时以长度平方为单位调整时间。Huisken 的加权面积把这种缩放与面积耗散结合起来。\footnote{Gerhard Huisken（1958--），德国数学家，研究几何分析、平均曲率流与数学相对论；生年见\href{https://www.ae-info.org/ae/Member/Huisken_Gerhard/CV}{Academia Europaea 履历}。}其导数是法向误差的负平方：权不再下降时，形状按同一比例收缩。单调量帮助识别极限模型，而弱演化还须说明保留的是零水平集、带重数的面积测度，还是更强的参数结构。这些选择决定了“同一个初值”和“唯一解”的具体含义。

\ProseParagraphBreak
本章使用\secref{g3:sec:0401}的子流形基本方程、\secref{g3:sec:1901}的面积变分，以及\secref{g3:sec:2601}和\secref{g3:sec:2602}的最大值原理与拟线性短时理论。末节还需\secref{g3:sec:2304}的变分流形及广义平均曲率。前三节的核心结论在正文证明；弱解的全局构造与粘性比较定理在相应陈述处给出具体来源。

% !TeX root = ../../Book3.tex
% 纲要标题：### 28.1 方程与光滑短时理论
% 纲要位置：几何学计划纲要.md，第 3285 行。
% 正文状态：初稿；尚未正式审核。

% 纲要内容开始
% 
% * 参数退化及法向速度
% * 图像方程的抛物性
% * 规范与短时存在
% * 面积下降及边界条件
% 
% 纲要内容结束

\section{方程与光滑短时理论}
\label{g3:sec:2801}

面积的一阶变分指出，平均曲率向量给出面积下降的法向方向。让浸入沿这个方向移动，所得方程却没有在全部分量上直接呈现强抛物性：沿切向移动参数，不改变眼前的子流形。解决办法是先规定如何表示附近的子流形，再求它的法向运动。本节用初始浸入的法丛作为固定未知量空间，使\cref{g3:thm:quasilinear-parabolic-short-time}能够直接应用。

\subsection{参数退化及法向速度}
\label{g3:subsec:2801:01}
设 \(F:M^m\to(N^{m+q},\bar g)\) 为光滑浸入，\(m\ge1\)，诱导度量为 \(g=F^*\bar g\)。沿用\secref{g3:sec:0401}的第二基本形式与平均曲率约定：
\[
B(X,Y)=\bigl(\bar\nabla_X\mathrm dF(Y)\bigr)^\perp,
\qquad H=\frac1m\operatorname{tr}_g B.
\]
几何流文献常把未除以维数的迹称为平均曲率向量。为了保留本书已有的 \(H\)，本章另记
\begin{equation}\label{g3:eq:mcf-trace-convention}
	\mathcal H=\operatorname{tr}_gB=mH.
\end{equation}
这使闭源上的面积一阶变分写成
\(\mathcal A'(V)=-\int_M\langle\mathcal H,V\rangle\,\mu_g\)。于是面积相对于法向场的 \(L^2\) 度量的负梯度正是 \(\mathcal H\)。

\begin{definition}\label{g3:def:mean-curvature-flow}
	设 \(M^m\) 为光滑无边界流形，\(m\ge1\)，\((N,\bar g)\) 为光滑无边界黎曼流形。光滑浸入族 \(F:M\times[0,T)\to N\) 若满足
	\begin{equation}\label{g3:eq:mean-curvature-flow}
		\partial_tF=\mathcal H_F,\qquad F(\cdot,0)=F_0,
	\end{equation}
	称为以 \(F_0\) 为初值的\term[pingjun qulv liu]{平均曲率流}[mean curvature flow]。若只规定 \((\partial_tF)^\perp=\mathcal H_F\)，则称该浸入族具有平均曲率法向速度；它还可以带有切向速度。
\end{definition}
\symidx{\(\mathcal H\)：未除以维数的平均曲率向量}

在 Euclid 环境中，\cref{g3:prop:gauss-decomposition}的 Gauss 公式给出
\begin{equation}\label{g3:eq:mcf-euclidean-coordinate}
	\mathcal H=g^{ij}\bigl(F_{ij}-\Gamma^k_{ij}(g)F_k\bigr)
	=-\Delta_gF.
\end{equation}
因此方程形式上是 \(F_t+\Delta_gF=0\)，但度量 \(g\) 由 \(F\) 自身诱导。对一般环境，在右端加入
\(g^{ij}\bar\Gamma^\alpha_{\beta\gamma}(F)F_i^\beta F_j^\gamma\)；环境联络项只含 \(F\) 的一阶导数。

\ProseParagraphBreak
记 \(P_F^\perp\) 为到浸入法空间的正交投影。曲率算子的二阶部分为 \(g^{ij}P_F^\perp\partial_i\partial_jF\)。在固定浸入处线性化，并把 \(\partial_i\partial_j\) 换成 \(-\xi_i\xi_j\)，得到主符号
\[
\sigma_2(\mathrm D\mathcal H_F)(\xi)V
=-|\xi|_g^2P_F^\perp V.
\]
它在法向量上给出热方程的符号，在切向量上为零。这个核正对应重新参数化的自由。

\begin{lemma}\label{g3:lem:mcf-reparametrization}
	设 \(M\) 闭，\(F_t:M\to N\) 是光滑浸入族，且 \((\partial_tF)^\perp=\mathcal H_F\)。唯一写成
	\[
	\partial_tF=\mathcal H_F+\mathrm dF(W_t),
	\]
	其中 \(W_t\) 是 \(M\) 上的光滑时变向量场。若 \(\phi_t\) 解
	\begin{equation}\label{g3:eq:mcf-tangential-ode}
		\partial_t\phi_t=-W_t\circ\phi_t,\qquad\phi_0=\operatorname{id}_M,
	\end{equation}
	则 \(\widetilde F_t=F_t\circ\phi_t\) 满足 \(\partial_t\widetilde F=\mathcal H_{\widetilde F}\)。在 \(F\) 的每个紧时间子区间上，\(\phi_t\) 都存在并且是微分同胚。
\end{lemma}
\begin{proof}
	\(\mathrm dF\) 在每点把 \(\mathrm TM\) 同构到浸入切空间，所以 \(W_t\) 存在、唯一且光滑。闭性使光滑时变向量场在有限时间内不会逃离流形；常微分方程的局部存在唯一性可沿有限时间区间延续，其逆由反向时间方程给出。链式法则与 \(\mathcal H_{F\circ\phi}=\mathcal H_F\circ\phi\) 给出
	\[
	\partial_t(F_t\circ\phi_t)
	=(\mathcal H_F+\mathrm dF(W_t))\circ\phi_t
	-\mathrm dF_{\phi_t}(W_t\circ\phi_t)
	=\mathcal H_{F_t\circ\phi_t}.
	\]
\end{proof}

时间不变的重新参数化保持式\eqref{g3:eq:mean-curvature-flow}；时间变化的重新参数化通常增加一个切向速度。给定初始映射后的参数化唯一性与不指定参数的几何运动，需要分别核对。

\subsection{图像方程的抛物性}
\label{g3:subsec:2801:02}
先在 Euclid 空间中观察一个可直接计算的规范。设 \(u:U\times[0,T)\to\mathbb R^q\)，其中 \(U\subset\mathbb R^m\) 开，令 \(F(x,t)=(x,u(x,t))\)。底部坐标始终固定，故速度是竖直的，而非通常意义下纯法向的。

\begin{proposition}\label{g3:prop:mcf-graph-equation}
	上述光滑图族具有平均曲率法向速度，当且仅当
	\begin{equation}\label{g3:eq:mcf-graph-system}
		\partial_tu^\alpha=g^{ij}(\mathrm Du)u^\alpha_{ij},
		\qquad g_{ij}=\delta_{ij}+\sum_{\alpha=1}^q u_i^\alpha u_j^\alpha,
		\quad 1\le\alpha\le q.
	\end{equation}
	它是任意余维下的拟线性强抛物系统；在 \(|\mathrm Du|\le K\) 的区域，主系数满足
	\[
	\frac{|\xi|^2}{1+K^2}|v|^2
	\le g^{ij}\xi_i\xi_j|v|^2\le |\xi|^2|v|^2
	\qquad(\xi\in\mathbb R^m,\ v\in\mathbb R^q).
	\]
\end{proposition}
\begin{proof}
	图的 \(F_{ij}=(0,u_{ij})\)，故
	\(\mathcal H=P_F^\perp(0,g^{ij}u_{ij})\)，而
	\((\partial_tF)^\perp=P_F^\perp(0,u_t)\)。竖直向量上的 \(P_F^\perp\) 是单射：若 \((0,a)\) 切于图，必有 \((0,a)=(X,\mathrm Du(X))\)，从而 \(X=0\)、\(a=0\)。两边的法向投影相等遂等价于式\eqref{g3:eq:mcf-graph-system}。

	对任意 \(X\in\mathbb R^m\)，有
	\(|X|^2\le g(X,X)\le(1+K^2)|X|^2\)。取逆矩阵就得到所述符号界；它在每个竖直分量上乘同一个正标量。故即使 \(q>1\)，这一图规范的最高阶部分仍是标量系数乘恒等算子。
\end{proof}

\begin{example}\label{g3:exm:mcf-scalar-graph}
	余维为一时，逆矩阵可以直接写出：
	\[
	g^{ij}=\delta^{ij}-\frac{u_i u_j}{1+|\mathrm Du|^2},\qquad
	u_t=\sqrt{1+|\mathrm Du|^2}\,
	\operatorname{div}\frac{\mathrm Du}{\sqrt{1+|\mathrm Du|^2}}.
	\]
	向上单位法向为 \(\nu=(-\mathrm Du,1)/\sqrt{1+|\mathrm Du|^2}\)。竖直速度的法向分量是
	\(u_t\nu/\sqrt{1+|\mathrm Du|^2}\)，恰等于 \(\mathcal H\)。若 \(u=a\cdot x+b\)，则 \(u_{ij}=0\)，图是静止的平面。主符号的最小特征值为 \((1+|\mathrm Du|^2)^{-1}\)，说明梯度界在统一抛物性中的作用。
\end{example}

局部图方程展示了退化如何消失。要得到一个闭浸入的解，还需把所有图用同一个全局未知量联系起来；仅在各坐标片分别求解，无法预先保证它们沿时间相容。初始浸入的法丛提供了这种统一表示。

\subsection{规范与短时存在}
\label{g3:subsec:2801:03}
固定光滑浸入 \(F_0:M^m\to(N^{m+q},\bar g)\)，其中 \(M\) 闭、\(q\ge1\)。令 \(E=\nu F_0\) 为它的法丛，配上由环境度量及联络诱导的丛度量与法联络。由于 \(F_0(M)\) 紧，存在 \(E\) 零截面的统一邻域，使
\[
\mathscr E(x,z)=\exp_{F_0(x)}^{\bar g}(z)
\]
光滑有定义。对足够小的光滑截面 \(u\)，记
\(F_u(x)=\mathscr E(x,u(x))\)，称它为相对于 \(F_0\) 的\term[fa tu]{法图}[normal graph]。小 \(C^1\) 邻域内的 \(F_u\) 仍为浸入。这里不要求环境完备：指数映射只需在紧像附近的短向量上定义。

\begin{lemma}\label{g3:lem:mcf-normal-graph-gauge}
	对上述闭浸入 \(F_0\)，每个充分 \(C^1\) 接近 \(F_0\) 的光滑映射 \(G:M\to N\) 都唯一表示成
	\begin{equation}\label{g3:eq:mcf-normal-graph-factorization}
		G=F_u\circ\chi,
	\end{equation}
	其中 \(u\) 是充分小的光滑法截面，\(\chi\) 是充分 \(C^1\) 接近恒等映射的微分同胚。对光滑有限维参数族 \(G\)，这两个对象也联合光滑依赖参数。
\end{lemma}
\begin{proof}
	在 \((x,0)\in E\) 处，\(\mathscr E\) 的微分为
	\[
	\mathrm D\mathscr E_{(x,0)}(X,z)=\mathrm dF_0(X)+z.
	\]
	切空间与法空间的直和使它成为到 \(\mathrm T_{F_0(x)}N\) 的线性同构。有限维逆函数定理因而在每点给出一个局部逆。用 \(M\) 的有限覆盖缩小邻域后，对每个接近 \(F_0\) 的 \(G\)，可以唯一求得靠近 \((x,0)\) 的 \((y(x),z(x))\)，使
	\(G(x)=\mathscr E(y(x),z(x))\)。这些局部解在重叠处一致，因为都在靠近参数对角线的同一逆函数邻域内。具体地，先取较小的有限坐标片，使其闭包落在原逆函数片内；相邻片上的两个候选基点均充分靠近 \(x\)，从而同时属于一个原片，那里有局部逆的唯一性。于是 \(y,z\) 全局光滑，且 \(y\) 在 \(C^1\) 中接近恒等映射。

	充分 \(C^1\) 接近恒等映射的 \(y:M\to M\) 是微分同胚。局部微分可逆使它局部单射；用有限坐标片和逆函数估计，可以在靠近对角线的一致邻域内保证单射。对远离对角线的点对，\(y\) 与恒等映射的一致接近也保证像不重合。所以 \(y\) 全局单射，其像在每个分支中开且闭，故满射。令 \(\chi=y\)、\(u=z\circ\chi^{-1}\)，便得到式\eqref{g3:eq:mcf-normal-graph-factorization}。同一局部逆证明唯一性和参数光滑性。

	即使 \(F_0\) 有自交，以上论证仍成立：所选逆的基点始终靠近原参数 \(x\)。不同参数处的像可以重合，但并未要求 \(\mathscr E\) 在整个法丛邻域上单射。
\end{proof}

现在把法向速度方程写在固定丛 \(E\) 上。令
\[
L_u(x)=\mathrm D_z\mathscr E(x,u(x)):E_x\longrightarrow\mathrm T_{F_u(x)}N,
\qquad J_u=P_{F_u}^\perp L_u.
\]
在 \(u=0\) 处，\(J_0\) 就是 \(E_x\) 到原法空间的恒等同构；因而 \(u\) 的 \(C^1\) 范数足够小时，\(J_u:E_x\to\nu_xF_u\) 仍可逆。

\begin{lemma}\label{g3:lem:mcf-normal-graph-parabolic}
	在上述小法图邻域内，\((\partial_tF_u)^\perp=\mathcal H_{F_u}\) 等价于固定丛 \(E\) 上的方程
	\begin{equation}\label{g3:eq:mcf-normal-graph-pde}
		u_t=J_u^{-1}\mathcal H_{F_u}
		=g_u^{ij}\nabla_i\nabla_j u+\mathcal B(x,u,\nabla u),
		\qquad g_u=F_u^*\bar g,
	\end{equation}
	其中 \(\nabla\) 是固定的法联络，二阶导数同时使用 \(g_0=F_0^*\bar g\) 的源联络；\(\mathcal B\) 光滑依赖一阶喷射。此方程在零截面的一阶喷射邻域中一致强抛物。
\end{lemma}
\begin{proof}
	时间微分给出 \(\partial_tF_u=L_u u_t\)，故法向速度就是 \(J_u u_t\)。在环境局部坐标和 \(E\) 的固定局部标架中，链式法则把 \(F_u\) 的协变二阶导数写成
	\[
	\bar\nabla_i\partial_jF_u
	=L_u(\nabla_i\nabla_j u)+R_{ij}(x,u,\nabla u),
	\]
	其中固定源联络、丛联络及环境联络造成的差都只进入光滑低阶项 \(R_{ij}\)。取 \(g_u\) 的迹并投影到 \(F_u\) 的法空间，得到
	\[
	\mathcal H_{F_u}=g_u^{ij}J_u(\nabla_i\nabla_j u)
	+g_u^{ij}P_{F_u}^\perp R_{ij}.
	\]
	左乘 \(J_u^{-1}\) 即得方程。\(g_u\)、\(J_u\) 及其逆都只依赖一阶喷射，所以低阶项具有所述形式。取固定背景度量 \(g_0\)；紧性和小 \(C^1\) 邻域给出 \(c^{-1}g_0\le g_u\le c g_0\)。主符号因此为 \(g_u^{ij}\xi_i\xi_j\operatorname{Id}_E\)，具有一致正下界。这正是\cref{g3:def:quasilinear-strong-parabolicity}的条件。
\end{proof}

\begin{theorem}\label{g3:thm:mcf-short-time}
	设 \(M^m\) 为闭光滑流形，\(m\ge1\)，\((N^{m+q},\bar g)\) 为光滑无边界黎曼流形，\(q\ge1\)，\(F_0:M\to N\) 为光滑浸入。则存在 \(T>0\) 及光滑浸入族 \(F:M\times[0,T]\to N\)，满足
	\(\partial_tF=\mathcal H_F\)、\(F(\cdot,0)=F_0\)。具有相同初始映射的光滑解在共同存在区间内相同。若 \(F_0\) 为嵌入，则时间足够小时每个 \(F_t\) 也是嵌入。
\end{theorem}
\begin{proof}
	对式\eqref{g3:eq:mcf-normal-graph-pde}取初值 \(u(0)=0\)。未知量是固定正秩丛 \(E\) 的截面，主系数与低阶项在一个一阶喷射开邻域上光滑，并在初始喷射图上一致强抛物；闭性使初始喷射图紧。\Cref{g3:thm:quasilinear-parabolic-short-time}因此给出联合光滑的短时解。缩短时间可使其留在所选小 \(C^1\) 邻域，所以 \(\widehat F=F_u\) 始终是浸入，且具有平均曲率法向速度。

	唯一写成
	\(\partial_t\widehat F=\mathcal H_{\widehat F}+\mathrm d\widehat F(W_t)\)。由\cref{g3:lem:mcf-reparametrization}，解式\eqref{g3:eq:mcf-tangential-ode}后，\(F=\widehat F\circ\phi_t\) 满足所需纯法向方程。这一步是光滑常微分方程，保留已经得到的初时刻光滑性。

	为证唯一性，设 \(G\) 是同初值的另一光滑解。短时间内，\cref{g3:lem:mcf-normal-graph-gauge}给出 \(G=F_v\circ\chi_t\)，其中 \(v(0)=0\)、\(\chi_0=\operatorname{id}\)。微分这个分解并取法向部分，\(\chi_t\) 的时间导数只产生切向项，所以
	\(J_v v_t=\mathcal H_{F_v}\)。于是 \(v\) 满足与 \(u\) 相同的强抛物初值问题；\cref{g3:prop:quasilinear-parabolic-uniqueness}给出 \(v=u\)。再比较切向部分，\(G_t\) 纯法向迫使
	\(\partial_t\chi_t=-W_t\circ\chi_t\)。常微分方程唯一性给 \(\chi_t=\phi_t\)，从而 \(G=F\)。在共同存在区间内的任意两解相等时刻，都可用该时刻的浸入重做此局部论证；连续性和向右延续遂给出整个共同区间的唯一性。

	最后，闭流形的嵌入在 \(C^1\) 拓扑下是开条件。其原因仍是两类点对的分离：近点由初始嵌入的有限局部图与微分下界控制；离开对角线的点对，在紧集上有正的像间距。充分小的 \(C^1\) 扰动同时保持二者。因此初始嵌入时缩短 \(T\) 即得末项。
\end{proof}

当 \(q=0\) 时，法空间为零，\(\mathcal H\equiv0\)，式\eqref{g3:eq:mean-curvature-flow}的唯一解就是 \(F_t=F_0\)。正余维证明中的两个选择各有作用：法图使方程强抛物，随后源上的微分同胚使运动速度完全法向。自交的浸入同样适用，因为这些选择使用参数空间上的法丛。

\subsection{面积下降及边界条件}
\label{g3:subsec:2801:04}
短时解存在以后，面积下降才成为实际演化的恒等式。用 \(V=\partial_tF\) 表示一般浸入变分的速度，写 \(V=V^\perp+\mathrm dF(W)\)。由\cref{g3:eq:submanifold-area-first-variation}的点态计算，
\begin{equation}\label{g3:eq:mcf-general-area-density}
	\partial_t\mu_g
	=\bigl(\operatorname{div}_gW-\langle\mathcal H,V^\perp\rangle\bigr)\mu_g.
\end{equation}
散度项描述参数运动带来的密度输运；法向项描述真实的面积变化。

\begin{proposition}\label{g3:prop:mcf-area-dissipation}
	设 \(M^m\) 闭，\(F:M\times[0,T)\to(N,\bar g)\) 为光滑平均曲率流。令 \(g_t=F_t^*\bar g\)、\(\mathcal A(t)=\int_M\mu_{g_t}\)。则
	\begin{align}
		\partial_tg_{ij}&=-2\langle\mathcal H,B_{ij}\rangle,
		\label{g3:eq:mcf-metric-evolution}\\
		\partial_t\mu_{g_t}&=-|\mathcal H|^2\mu_{g_t},\qquad
		\mathcal A'(t)=-\int_M|\mathcal H|^2\,\mu_{g_t}.
		\label{g3:eq:mcf-area-dissipation}
	\end{align}
	因而对 \(0\le a<b<T\)，
	\[
	\mathcal A(a)-\mathcal A(b)
	=\int_a^b\int_M|\mathcal H|^2\,\mu_{g_t}\,\mathrm dt.
	\]
	面积在一个时间区间上恒定，当且仅当该区间上的流为静止极小浸入。
\end{proposition}
\begin{proof}
	微分诱导度量得到
	\(\partial_tg_{ij}=\langle\bar\nabla_i\mathcal H,F_j\rangle+
	\langle F_i,\bar\nabla_j\mathcal H\rangle\)。对 \(\langle\mathcal H,F_j\rangle=0\) 求导，并用 \(B\) 的对称性，便得第一式。再取一半的迹，得到
	\(\partial_t\mu_g=\tfrac12g^{ij}\partial_tg_{ij}\mu_g=-|\mathcal H|^2\mu_g\)。闭流形上积分与时间微分可交换，所得等式再对时间积分即为耗散恒等式。若面积恒定，连续非负函数 \(|\mathcal H|^2\) 的积分在每时刻为零，故 \(\mathcal H\equiv0\)，方程遂给 \(\partial_tF=0\)；逆向立即成立。
\end{proof}

面积控制的是 \(\mathcal H\) 的时空平方积分。它不能单独给出第二基本形式的点态上界：某些方向的弯曲可以在取迹时抵消，曲率也可能集中在越来越小的区域。后续演化公式与加权面积分别提供更局部的观察。

\ProseParagraphBreak
若 \(M\) 紧且有光滑边界，对具有平均曲率法向速度的光滑浸入族积分式\eqref{g3:eq:mcf-general-area-density}，得到
\begin{equation}\label{g3:eq:mcf-boundary-flux}
	\frac{\mathrm d}{\mathrm dt}\mathcal A(t)
	=-\int_M|\mathcal H|^2\,\mu_g
	+\int_{\partial M}\langle V,\mathrm dF(\eta)\rangle\,\mu_{\partial g},
\end{equation}
其中 \(\eta\) 是 \(M\) 内的外单位余法向。固定边界 \(F_t|_{\partial M}=F_0|_{\partial M}\) 使 \(V=0\)，通量消失。若还要求纯法向方程光滑成立到边界，则必有 \(\mathcal H=0\) 于边界；这已经是初始数据需要满足的相容条件之一。

\ProseParagraphBreak
另一种情形是边界留在固定光滑支撑超曲面 \(S\subset N\)，并满足正交接触条件 \(\mathrm dF(\eta)\perp\mathrm TS\)。边界速度属于 \(\mathrm TS\)，所以式\eqref{g3:eq:mcf-boundary-flux}的边界项仍为零。这就是\term[ziyou bianjie pingjun qulv liu]{自由边界平均曲率流}[free boundary mean curvature flow]中面积耗散的来源。这两种边界情形的公式以已有光滑解为前提；它们的短时构造需要配套的抛物边值理论和初始相容条件，而\cref{g3:thm:mcf-short-time}处理闭源。

\begin{example}\label{g3:exm:mcf-free-boundary-flux}
	上半球面 \(F(x,t)=R(t)x\)，\(x\in S^m\)、\(x_{m+1}\ge0\)，以赤道为边界，支撑面为 \(S=\{x_{m+1}=0\}\)。取 \(R'(t)=-m/R(t)\)，则速度等于 \(\mathcal H=-mx/R(t)\)。在赤道，球面的外余法向为 \(-e_{m+1}\)，它垂直于支撑面；边界径向速度则切于该平面。因此通量为零，半球面积满足同一耗散恒等式。反之，\cref{g3:exm:plane-disk-moving-boundary}中的扩张平盘具有零平均曲率法向速度，却因正的边界通量而增加面积。这两个计算显示，边界怎样移动是耗散公式的一部分。
\end{example}
% 纲要条目排为编号小节标题，并进入目录与书签；说明段仍保留于纲要注释中。

% !TeX root = ../../Book3.tex
% 纲要标题：### 28.2 演化公式与比较
% 纲要位置：几何学计划纲要.md，第 3292 行。
% 正文状态：初稿；尚未正式审核。

% 纲要内容开始
% 
% * 第二基本形式和平均曲率的演化
% * 超曲面避免相交原理
% * 凸性与平均凸性的保持
% * 曲线缩短流和球面模型
% 
% 纲要内容结束

\section{演化公式与比较}
\label{g3:sec:2802}

面积下降只给出总量的控制。一个很小的区域可以弯得越来越急，同时总面积继续减小；要判断这种局部变化，需要把方程作用到第二基本形式上。本节的演化与比较均在 Euclid 环境中进行。对超曲面，法空间只有一个方向，向量方程可以化成带符号的标量曲率方程；两个超曲面在距离极小处的切空间又恰好相同，这使曲率计算能够转成避免相交的比较结论。

\subsection{第二基本形式和平均曲率的演化}
\label{g3:subsec:2802:01}
设 \(F:M^m\times[0,T)\to\mathbb R^{m+1}\) 是光滑浸入的平均曲率流，\(m\ge1\)，并在所讨论的区域选取随时间光滑变化的单位法向 \(\nu\)。记
\[
	h_{ij}=\langle B(\partial_i,\partial_j),\nu\rangle,
	\qquad A^i{}_j=g^{ik}h_{kj},
	\qquad s=\operatorname{tr}A,
	\qquad q=|A|^2=\operatorname{tr}(A^2).
\]
于是 \(\mathcal H=s\nu\)、\(\partial_tF=s\nu\)，而\secref{g3:sec:1901}中的归一化平均曲率为 \(H=(s/m)\nu\)。形算子仍取 \(A=-\mathrm d\nu\)；外向半径 \(R\) 球面上 \(A=-R^{-1}\operatorname{Id}\)、\(s=-m/R\)。以下 \(\Delta=\nabla^*\nabla\) 表示当前度量的非负粗 Laplace 算子，在函数上就是 \(-\operatorname{div}\nabla\)。各个时刻的联络、范数与迹均使用 \(g(t)\)。

\begin{lemma}\label{g3:lem:mcf-simons-identity}
	设 \(F:M^m\to\mathbb R^{m+1}\) 为光滑双侧浸入，\(m\ge1\)，选单位法向 \(\nu\)，记 \(g=F^*g_{\mathrm E}\)、\(h_{ij}=\langle B_{ij},\nu\rangle\)、\(A^i{}_j=g^{ik}h_{kj}\)、\(s=\operatorname{tr}A\)、\(q=|A|^2\)，\(\Delta=\nabla^*\nabla\) 为非负粗 Laplace 算子。则第二基本形式满足
	\begin{equation}\label{g3:eq:mcf-simons-identity}
		-\Delta h_{ij}=\nabla_i\nabla_j s+s h_i{}^k h_{kj}-q h_{ij}.
	\end{equation}
\end{lemma}
\begin{proof}
	由\cref{g3:thm:codazzi-equation}，Euclid 环境中的 \(\nabla h\) 对三个输入对称，因而 \(\nabla^k h_{kj}=\nabla_j s\)。在所算点取正交标架，并令联络系数在该点为零。交换导数给出
	\[
		\sum_k\nabla_k\nabla_k h_{ij}
		=\nabla_i\nabla_j s+
		\sum_k([\nabla_k,\nabla_i]h)_{kj}.
	\]
	协变二阶张量的交换项是
	\[
		([\nabla_k,\nabla_i]h)_{kj}
		=-h(R(e_k,e_i)e_k,e_j)-h(e_k,R(e_k,e_i)e_j).
	\]
	\cref{g3:thm:gauss-equation}在当前余维与环境下写成
	\(R(X,Y)Z=h(Y,Z)AX-h(X,Z)AY\)。代入后，第一组收缩为
	\(- (h^3)_{ij}+s(h^2)_{ij}\)，第二组为
	\(-q h_{ij}+(h^3)_{ij}\)，其中乘积均用 \(g\) 收缩。两个三次矩阵乘积抵消，所余正是式\eqref{g3:eq:mcf-simons-identity}的反应项。左端的导数迹等于 \(-\Delta h\)，故结论成立。
\end{proof}

这个恒等式把平均曲率的 Hessian 改写成整个第二基本形式的扩散项。时间变分首先产生前者，而抛物估计需要后者。

\begin{proposition}\label{g3:prop:mcf-hypersurface-evolution}
	设 \(F:M^m\times[0,T)\to\mathbb R^{m+1}\) 为光滑双侧浸入的平均曲率流，\(m\ge1\)，满足 \(\partial_tF=s\nu\)。记 \(g=F_t^*g_{\mathrm E}\)、\(h_{ij}=\langle B_{ij},\nu\rangle\)、\(A=g^{-1}h\)、\(s=\operatorname{tr}A\)、\(q=|A|^2\)，\(\Delta=\nabla^*\nabla\) 为当前非负粗 Laplace。在固定源坐标中有
	\begin{align}
		\partial_t g_{ij}&=-2s h_{ij},&
		\partial_t g^{ij}&=2s h^{ij},&
		\partial_t\nu&=-\nabla s,\label{g3:eq:mcf-metric-normal-evolution}\\
		\partial_t h_{ij}&=\nabla_i\nabla_j s-s h_i{}^k h_{kj},&
		\partial_t\mu_g&=-s^2\mu_g.\label{g3:eq:mcf-second-form-variation}
	\end{align}
	相应的抛物演化为
	\begin{align}
		(\partial_t+\Delta)h_{ij}&=q h_{ij}-2s h_i{}^k h_{kj},\label{g3:eq:mcf-covariant-form-evolution}\\
		(\partial_t+\Delta)A^i{}_j&=q A^i{}_j,\label{g3:eq:mcf-shape-evolution}\\
		(\partial_t+\Delta)s&=q s,\label{g3:eq:mcf-mean-evolution}\\
		(\partial_t+\Delta)q&=-2|\nabla A|^2+2q^2.\label{g3:eq:mcf-curvature-norm-evolution}
	\end{align}
	这些式子均为内点恒等式，不要求 \(M\) 紧。
\end{proposition}
\begin{proof}
	写 \(F_i=\partial_iF\)。由 \(\partial_i\nu=-h_i{}^kF_k\)，对 \(g_{ij}=\langle F_i,F_j\rangle\) 求导即得第一式。对 \(g^{ik}g_{kj}=\delta^i_j\) 求导得第二式。再对 \(\langle\nu,F_i\rangle=0\)、\(\langle\nu,\nu\rangle=1\) 求导，得到 \(\langle\partial_t\nu,F_i\rangle=-\partial_i s\) 及 \(\partial_t\nu\perp\nu\)，故 \(\partial_t\nu=-\nabla s\)。

	在一点采用源的法坐标。对 \(h_{ij}=\langle\partial_i\partial_jF,\nu\rangle\) 求时间导数，其中 \(\partial_i\partial_j(s\nu)\) 的法向分量为
	\(\partial_i\partial_j s+s\langle\partial_i\partial_j\nu,\nu\rangle\)。由 \(|\nu|=1\)，后一个内积等于 \(-\langle\partial_i\nu,\partial_j\nu\rangle=-h_i{}^k h_{kj}\)。另一项 \(\langle\partial_i\partial_jF,-\nabla s\rangle\) 在该点为零；一般坐标下它正好把普通 Hessian 补成协变 Hessian。这证明 \(h\) 的变分式。面积元公式则由
	\(\partial_t\mu_g=\frac12g^{ij}\partial_tg_{ij}\mu_g\) 得出。

	把\cref{g3:lem:mcf-simons-identity}代入 \(h\) 的变分式，得到式\eqref{g3:eq:mcf-covariant-form-evolution}。对 \(A^i{}_j=g^{ik}h_{kj}\) 求导时，逆度量贡献 \(2s(A^2)^i{}_j\)，恰抵消该式中的最后一项。因为 \(\nabla g=0\)，升指标与空间 Laplace 运算相容，于是得到式\eqref{g3:eq:mcf-shape-evolution}。取混合张量的迹便得 \(s\) 的方程。

	最后用 \(q=\operatorname{tr}(A^2)\) 求导。对当前度量的空间协变导数，乘积法则给出
	\[
		-\Delta q=-2\operatorname{tr}(A\Delta A)+2|\nabla A|^2.
	\]
	另一方面，\(\partial_tq=2\operatorname{tr}(A\partial_tA)\)。将 \(A\) 的方程代入，得到式\eqref{g3:eq:mcf-curvature-norm-evolution}。这里使用混合张量的迹，已包含范数随度量变化的贡献。
\end{proof}

标量 \(s\) 可以有正有负，而 \(q\) 非负。其演化式中的 \(-2|\nabla A|^2\) 有利于压低空间振荡，\(2q^2\) 却使曲率可能增大；两项竞争与面积总量的下降可以同时发生。对任意闭源上的光滑流，固定 \(T_1<T\)，\(g(t)\) 在 \([0,T_1]\) 上一致正定。用一个固定背景联络表达 \(\Delta_{g(t)}\)，所得时间依赖系数满足\cref{g3:thm:parabolic-maximum}的条件，因此该定理也能用于这里的标量量。

\begin{corollary}\label{g3:cor:mcf-curvature-ode-bound}
	设 \(M^m\) 闭，\(m\ge1\)，\(F:M\times[0,T)\to\mathbb R^{m+1}\) 是光滑双侧浸入的平均曲率流，\(Q_0=\max_M|A(0)|^2>0\)。则对 \(0\le t<\min\{T,(2Q_0)^{-1}\}\)，
	\[
		|A(t)|^2\le\frac{Q_0}{1-2Q_0t}.
	\]
\end{corollary}
\begin{proof}
	令 \(Q(t)=Q_0/(1-2Q_0t)\)，则 \(Q'=2Q^2\)。由式\eqref{g3:eq:mcf-curvature-norm-evolution}，\(w=q-Q\) 满足
	\[
		(\partial_t+\Delta)w\le 2(q+Q)w,\qquad w(0)\le0.
	\]
	在任一所述范围内的紧时间区间上，零阶系数有界，故标量最大值原理给出 \(w\le0\)。任意这样的区间都适用，结论成立。
\end{proof}

\subsection{超曲面避免相交原理}
\label{g3:subsec:2802:02}
如果两张紧超曲面最接近的两点仍有正距离，那么连接这两点的线段同时垂直于两个切空间。余维为一使这两个切空间相同。沿共同的切方向把两点一起移动，就能将距离的二阶变化与两张超曲面的平均曲率比较。

\begin{theorem}\label{g3:thm:mcf-avoidance}
	设 \(M^m,N^m\) 为闭光滑流形，\(m\ge1\)，\(F:M\times[0,T)\to\mathbb R^{m+1}\)、\(G:N\times[0,T)\to\mathbb R^{m+1}\) 为光滑浸入的平均曲率流。若 \(F(M,0)\cap G(N,0)=\varnothing\)，则
	\[
		d(t)=\min_{p\in M,\,q\in N}|F(p,t)-G(q,t)|
	\]
	在 \([0,T)\) 上为正且不减。特别地，两张闭嵌入超曲面在共同的光滑存在期间保持不相交。这一结论称为\term[bimian xiangjiao yuanli]{避免相交原理}[avoidance principle]。
\end{theorem}
\begin{proof}
	记 \(z=F(p,t)-G(q,t)\)、\(f(p,q,t)=|z|^2/2\)。紧性保证 \(f\) 在每个时刻有最小值。先限制在 \(d(t)>0\) 的时间区间。在一个空间最小点，\(\langle z,\mathrm dF(X)\rangle=\langle z,\mathrm dG(Y)\rangle=0\)，所以两张超曲面的切空间均为 \(z^\perp\)。在这个共同空间取正交单位基 \(e_1,\ldots,e_m\)，并在两个源切空间分别取其对应的单位向量 \(X_i,Y_i\)。

	在 \(M\times N\) 中沿 \((X_i,Y_i)\) 的乘积测地线求二阶导数。在所算点两个一阶速度的差为零，因此
	\[
		\operatorname{Hess}f\bigl((X_i,Y_i),(X_i,Y_i)\bigr)
		=\langle z,B_F(X_i,X_i)-B_G(Y_i,Y_i)\rangle.
	\]
	最小点的 Hessian 非负，求和便得
	\[
		0\le\langle z,\mathcal H_F-\mathcal H_G\rangle
		=\partial_t f.
	\]
	为直接从这个点态结论推出最小值不减，取上述时间区间内 \(a<b\)。若 \(\min f(b)<\min f(a)\)，可选 \(\varepsilon>0\)，使 \(\min f(b)+\varepsilon(b-a)<\min f(a)\)。函数 \(f+\varepsilon(t-a)\) 在 \(M\times N\times[a,b]\) 的最小值便在时间 \(t>a\) 取得；在那里空间 Hessian 非负、时间左导数非正，于是 \(\partial_tf\le-\varepsilon\)。这与上式矛盾。因此最小值不减。

	初始两个像都是紧集，故 \(d(0)>0\)。若存在第一次距离为零的时刻，在它之前刚证的单调性给出 \(d(t)\ge d(0)\)；连续性使该时刻仍有 \(d\ge d(0)>0\)，矛盾。于是正距离与单调性在全部共同存在时间成立。
\end{proof}

证明允许每张浸入自身有自交，但使用了两张像之间的正距离。它也说明余维假设的位置：高余维时，两切空间虽然都垂直于 \(z\)，却未必相同；共同移动基点时，两个切速度的差项就不再为零。带边界问题还可能在边界取得最短距离，因而需要另加边界分离或适当的边界比较条件。

\begin{example}\label{g3:exm:mcf-higher-codimension-intersection}
	取 \(a>0\)、\(R_0>a/2\)，在 \(\mathbb R^3\) 中令 \(R(t)=\sqrt{R_0^2-2t}\)，并考虑两个圆
	\[
		F(u,t)=(R(t)\cos u,R(t)\sin u,0),\qquad
		G(v,t)=(a+R(t)\cos v,0,R(t)\sin v).
	\]
	它们各自在固定平面内沿曲率向量缩小，因为 \(R'=-1/R\)。两平面的交集为 \(x\) 轴，两个圆与此轴的交点横坐标分别为 \(\{R,-R\}\) 和 \(\{a+R,a-R\}\)。初时刻这两组数互异，但在
	\[
		t_*=\frac12\left(R_0^2-\frac{a^2}{4}\right)
	\]
	有 \(R(t_*)=a/2\)，两圆在 \((a/2,0,0)\) 相交。此时半径仍为正，两条流都光滑。这给出两张闭子流形在余维二时违反避免相交结论的具体例子。
\end{example}

\begin{corollary}\label{g3:cor:mcf-preserves-embeddedness}
	设 \(M^m\) 闭，\(m\ge1\)，\(F:M\times[0,T)\to\mathbb R^{m+1}\) 是光滑浸入的平均曲率流。若 \(F(\cdot,0)\) 是嵌入，则每个 \(F(\cdot,t)\) 都是嵌入。
\end{corollary}
\begin{proof}
	固定 \(T_1<T\)，并在 \(M\) 上选一个固定背景距离 \(d_0\)。浸入性的局部坐标表示及 \(M\times[0,T_1]\) 的紧性给出 \(\delta,c>0\)，使同一连通分支上 \(0<d_0(p,q)\le\delta\) 时
	\[
		|F(p,t)-F(q,t)|\ge c\,d_0(p,q)
		\quad(0\le t\le T_1).
	\]
	确切地说，在每个点选择一个使 \(\mathrm dF\) 的某个 \(m\) 维坐标投影可逆的环境坐标块；缩小源坐标球后，投影导数与其中心值足够接近，积分线段上的导数即给出局部双 Lipschitz 下界。有限覆盖给出上述共同常数。

	现在在紧集 \(K=\{(p,q):d_0(p,q)\ge\delta\}\) 上比较 \(|F(p,t)-F(q,t)|^2/2\)；不同分支的所有点对一并计入 \(K\)。其空间边界上的距离至少为 \(c\delta\)，初时刻则由嵌入与紧性有正下界。选一个严格小于这两个下界的正数 \(b\)。若距离以后第一次降到 \(b\)，最小点必在 \(K\) 的内部，且两点间距仍为正。前一定理中共同切方向的二阶计算因此照用；给平方距离加上任意小的正时间斜率，就排除了低于初值与边界值的内部最小值。因此距离始终大于或等于 \(b\)。所有不同点对的像都不同，紧源上的单射浸入就是嵌入。\(T_1\) 任意，结论成立。
\end{proof}

\subsection{凸性与平均凸性的保持}
\label{g3:subsec:2802:03}
设一张双侧超曲面的单位法向已经选定。与外向凸球一致，令 \(k=-h\)、\(\theta=-s=\operatorname{tr}_g k\)。若 \(\theta\ge0\)，称它相对于该法向为\term[pingjun tuxing]{平均凸的}[mean convex]；若对所有切向量 \(v\) 都有 \(k(v,v)\ge0\)，则称其第二基本形式弱凸。严格不等号分别给出严格平均凸与严格凸。对于一个给定有界区域的光滑边界，使用外向单位法向。弱凸逐方向限制全部主曲率，平均凸则只限制它们的和。

\begin{theorem}\label{g3:thm:mcf-mean-convexity}
	设 \(M^m\) 闭，\(m\ge1\)，\(F:M\times[0,T)\to\mathbb R^{m+1}\) 为光滑双侧浸入的平均曲率流，并固定连续延拓的单位法向。记 \(s=\langle\mathcal H,\nu\rangle\)，其中 \(\mathcal H\) 为迹平均曲率向量。若 \(\theta(0)=-s(0)\ge0\)，则 \(\theta(t)\ge0\)。若初时刻处处 \(\theta(0)>0\)，则以后仍处处为正。
\end{theorem}
\begin{proof}
	由式\eqref{g3:eq:mcf-mean-evolution}，\((\partial_t+\Delta)\theta=q\theta\)。在任意 \([0,T_1]\)、\(T_1<T\) 上，\(q\) 有界，对 \(-\theta\) 应用\cref{g3:thm:parabolic-maximum}，得到第一项。若 \(a=\min_M\theta(0)>0\)，则 \(w=a-\theta\) 满足
	\[
		(\partial_t+\Delta)w=q w-q a\le q w,
		\qquad w(0)\le0.
	\]
	相同的比较给出 \(\theta\ge a>0\)。
\end{proof}

\begin{example}\label{g3:exm:mcf-mean-convex-torus}
	在 \(\mathbb R^3\) 中取 \(R>2r>0\) 的旋转环面
	\[
		F(u,v)=\bigl((R+r\cos v)\cos u,(R+r\cos v)\sin u,r\sin v\bigr).
	\]
	其外向法向为 \(\nu=(\cos v\cos u,\cos v\sin u,\sin v)\)。两个坐标方向正交，度量系数为 \((R+r\cos v)^2,r^2\)，第二基本形式系数为 \(-(R+r\cos v)\cos v,-r\)。于是
	\[
		\theta=\frac1r+\frac{\cos v}{R+r\cos v}
		\ge\frac1r-\frac1{R-r}>0.
	\]
	但在内圈 \(v=\pi\) 处，\(k\) 的两个主值为 \(1/r\) 与 \(-1/(R-r)\)，一正一负。它是严格平均凸的，却不弱凸。前一定理保证其光滑流期间保留平均凸性，并未把这个初始环面变成凸超曲面。
\end{example}

凸性的保持需要检查一个二次型在任意方向上的符号。把 \(h_{ij}\) 的各个坐标分量分别当作标量比较，既不保留坐标不变性，也不能控制最小特征值。下面直接在可能出现的第一个零方向上使用演化式。

\begin{theorem}\label{g3:thm:mcf-convexity}
	设 \(M^m\) 闭，\(m\ge1\)，\(F:M\times[0,T)\to\mathbb R^{m+1}\) 为光滑双侧浸入的平均曲率流，并按选定单位法向 \(\nu\) 令 \(h_{ij}=\langle B_{ij},\nu\rangle\)、\(k=-h\)。若 \(k(0)\) 为非负二次型，则 \(k(t)\) 在全部光滑存在时间内非负。若 \(k(0)\) 正定，则 \(k(t)\) 始终正定。
\end{theorem}
\begin{proof}
	记 \(L=-\Delta=\operatorname{tr}_g\nabla^2\)，由前面的公式，
	\[
		(\partial_t-L)k=q k-2\theta k^2,
		\qquad \partial_tg=-2\theta k,
	\]
	其中 \((k^2)_{ij}=k_i{}^\ell k_{\ell j}\)。固定 \(T_1<T\)，选 \(C>\max_{M\times[0,T_1]}q\)，并设
	\(S=k+a g\)、\(a(t)=\varepsilon e^{Ct}\)，\(\varepsilon>0\)。初时刻 \(S\) 正定。若它第一次在 \((x_*,t_*)\) 出现零方向，取 \(g(t_*)\) 单位向量 \(v\) 满足 \(S(v,v)=0\)。由于 \(S\) 在此时仍非负，\(S(v,w)=0\) 对所有 \(w\) 成立，因而
	\[
		k(v,w)=-a g(v,w),\qquad k^2(v,v)=a^2.
	\]
	将 \(v\) 在附近延拓成向量场 \(V\)，使 \(\nabla V=0\) 于所算点。标量 \(S(V,V)\) 此前非负，此时有零的空间最小值；其时间导数非正、\(L\) 非负。在 \((\partial_t-L)S(V,V)\) 的展开中，含 \(S(v,\cdot)\) 或 \(\nabla V\) 的附加项均为零。因此
	\[
		0\ge ((\partial_t-L)S)(v,v)
		=-qa-2\theta a^2+Ca+2\theta a^2
		=a(C-q)>0,
	\]
	矛盾。故 \(k+\varepsilon e^{Ct}g\) 始终正定；令 \(\varepsilon\downarrow0\)，得到 \(k\ge0\)。

	若 \(k(0)\) 正定，紧性给出 \(a_0>0\)，使 \(k(0)>a_0g(0)\)。改取 \(S=k-a(t)g\)、\(a(t)=a_0e^{-t}\)。若出现第一个零方向，则 \(k(v,w)=a g(v,w)\)、\(k^2(v,v)=a^2\)，同样计算得到
	\[
		0\ge ((\partial_t-L)S)(v,v)
		=qa-2\theta a^2-a'+2\theta a^2
		=a(q+1)>0.
	\]
	因此 \(k(t)>a_0e^{-t}g(t)\)，特别地仍正定。任意 \(T_1<T\) 都可如此处理。
\end{proof}

\subsection{曲线缩短流和球面模型}
\label{g3:subsec:2802:04}
当 \(m=1\) 时，平均曲率向量就是曲率向量，流也称为\term[quxian suoduan liu]{曲线缩短流}[curve shortening flow]。设 \(F:S^1\times[0,T)\to\mathbb R^2\) 是闭正则曲线的光滑流，取单位法向 \(\nu\) 并记带符号曲率为 \(\kappa\)，于是 \(\partial_tF=\kappa\nu\)。以 \(\sigma\) 表示当前弧长，\cref{g3:prop:mcf-hypersurface-evolution}化为
\begin{equation}\label{g3:eq:curve-shortening-evolution}
	\partial_t(\mathrm d\sigma)=-\kappa^2\,\mathrm d\sigma,
	\qquad \partial_t\kappa=\partial_\sigma^2\kappa+\kappa^3,
	\qquad \frac{\mathrm d}{\mathrm dt}\operatorname{Length}(F_t)
	=-\int_{S^1}\kappa^2\,\mathrm d\sigma.
\end{equation}
这里 \(\partial_t\) 固定原来的源参数。若该参数为 \(u\)，令 \(v=|F_u|\)，则 \(\partial_\sigma=v^{-1}\partial_u\)、\(v_t=-\kappa^2v\)，故
\[
	[\partial_t,\partial_\sigma]=\kappa^2\partial_\sigma.
\]
弧长参数随时间变化，在继续交换导数时要保留这一项。

\begin{example}\label{g3:exm:curve-shortening-graph-translator}
	对平面图像 \(F(x,t)=(x,u(x,t))\)，取向上单位法向
	\(\nu=(-u_x,1)/(1+u_x^2)^{1/2}\)。其曲率与图像的法向速度分别为
	\[
		\kappa=\frac{u_{xx}}{(1+u_x^2)^{3/2}},
		\qquad \langle\partial_tF,\nu\rangle=
		\frac{u_t}{(1+u_x^2)^{1/2}}.
	\]
	因此曲线缩短流的图像方程为
	\[
		u_t=\frac{u_{xx}}{1+u_x^2}.
	\]
	在 \(-\pi/2<x<\pi/2\) 上，\(u(x,t)=-\log\cos x+t\) 给出一个显式解：\(u_x=\tan x\)、\(u_{xx}=\sec^2x\)，右端恒为 \(1\)。其图像竖直平移而形状不变，是一个非紧平移解。固定图像参数带有切向速度；消去该速度后才满足 \(\partial_tF=\kappa\nu\) 的纯法向参数形式。该曲线长度无穷，闭曲线的总长度积分式不能直接用于它。
\end{example}

\begin{example}\label{g3:exm:mcf-shrinking-sphere}
	对 \(S^m\subset\mathbb R^{m+1}\)、\(m\ge1\)，令 \(F(\omega,t)=R(t)\omega\)，取外向法向 \(\nu=\omega\)。由 \(s=-m/R\)，流方程成为
	\[
		R'=-\frac mR,
		\qquad R(t)=\sqrt{R_0^2-2mt},
		\qquad 0\le t<T_*:=\frac{R_0^2}{2m}.
	\]
	面积和曲率分别为
	\[
		\mathcal A(t)=|S^m|(R_0^2-2mt)^{m/2},
		\qquad |A(t)|^2=\frac m{R_0^2-2mt}
		=\frac1{2(T_*-t)}.
	\]
	面积下降到零，曲率却发散；初始的严格凸性在 \(t<T_*\) 始终保留。直接求导也验证
	\(\mathcal A'=-m^2|S^m|R^{m-2}=-\int s^2\mu_g\)。当 \(m=1\) 时，这就是半径 \(\sqrt{R_0^2-2t}\) 的缩圆。
\end{example}

闭性本身已经使 Euclid 环境中的光滑平均曲率流具有有限的可能寿命。下面的计算不限余维，也不要求凸性。

\begin{proposition}\label{g3:prop:mcf-finite-lifetime}
	设 \(M^m\) 非空且闭，\(m\ge1\)，\(F:M\times[0,T)\to\mathbb R^{m+q}\)、\(q\ge1\)，为光滑浸入的平均曲率流。若初像包含于以 \(a\) 为中心、半径 \(R_0\) 的闭球，则
	\begin{equation}\label{g3:eq:mcf-shrinking-ball-bound}
		|F(x,t)-a|^2\le R_0^2-2mt
		\quad(0\le t<T),
	\end{equation}
	从而 \(T\le R_0^2/(2m)\)。
\end{proposition}
\begin{proof}
	Euclid 位置向量满足 \(-\Delta F=\mathcal H\)，并且 \(|\mathrm dF|_g^2=m\)。因此
	\[
		-\Delta|F-a|^2=2\langle F-a,\mathcal H\rangle+2m,
		\qquad \partial_t|F-a|^2=2\langle F-a,\mathcal H\rangle.
	\]
	相减得到 \((\partial_t+\Delta)(|F-a|^2+2mt)=0\)。在任意紧时间区间用标量最大值原理，就得式\eqref{g3:eq:mcf-shrinking-ball-bound}。左边非负，所以不能存在满足 \(t>R_0^2/(2m)\) 的流时刻。
\end{proof}

球面使上述寿命上界成为等号，也给出曲率变大的确切速率。对一般初始超曲面，保持平均凸或凸只是一种随时间不变的曲率条件；要进一步确定奇点形状，需要在缩小的空间尺度上比较整个流。下一节的加权面积会同时记录时间距离、空间位置与面积缩放。
% 纲要条目排为编号小节标题，并进入目录与书签；说明段仍保留于纲要注释中。

% !TeX root = ../../Book3.tex
% 纲要标题：### 28.3 Huisken 单调公式
% 纲要位置：几何学计划纲要.md，第 3299 行。
% 正文状态：初稿，尚未正式审核。

% 纲要内容开始
% 
% * 后向热核及加权面积
% * 微分和分部积分
% * 平方项与等号条件
% * 自收缩子和奇点放大
% 
% 纲要内容结束

\section{Huisken 单调公式}
\label{g3:sec:2803}

总面积下降记录整个子流形的变化。若曲率只在很小的区域集中，总面积却可能变化不大。为观察指定时空点附近的形状，需要随时间缩小观察范围，同时消去面积在缩放中的量纲。高斯权同时完成这两件事。本节固定 Euclid 环境 \(\mathbb R^N\)，流速为 \(\mathcal H=\operatorname{tr}_gB\)；闭流的积分均在参数流形上计算，因此也适用于有自交的浸入。

\subsection{后向热核及加权面积}
\label{g3:subsec:2803:01}
给定 \(x_0\in\mathbb R^N\) 与时间 \(t_0\)，置 \(\tau=t_0-t>0\)，定义
\begin{equation}\label{g3:eq:mcf-backward-gaussian}
	\rho_{x_0,t_0}(x,t)
	\coloneqq(4\pi\tau)^{-m/2}
	\exp\!\left(-\frac{|x-x_0|^2}{4\tau}\right).
\end{equation}
这里 \(m\) 是移动子流形的维数。此函数在经过 \(x_0\) 的任意 \(m\) 维仿射平面上积分为一；在整个 \(\mathbb R^N\) 上则通常不是归一化的 \(N\) 维热核。取 \(m/2\) 为指数，是为了使权与 \(m\) 维面积的缩放相抵消。

\begin{definition}\label{g3:def:mcf-gaussian-area}
	固定 \(x_0\in\mathbb R^N\)、\(t_0\in\mathbb R\)，记 \(\rho_{x_0,t_0}(y,t)=[4\pi(t_0-t)]^{-m/2}\exp(-|y-x_0|^2/[4(t_0-t)])\)，\(g(t)=F_t^*g_{\mathrm E}\)、\(\mathcal H=\operatorname{tr}_{g(t)}B\)。
	设 \(F:M^m\times[0,T)\to\mathbb R^N\) 是闭光滑流形的光滑平均曲率流。对 \(t<\min(T,t_0)\)，定义其\term[gaosi jiaquan mianji]{高斯加权面积}[Gaussian weighted area]为
	\[
		\Phi_{x_0,t_0}(t)
		\coloneqq\int_M\rho_{x_0,t_0}(F(p,t),t)\,\mu_{g(t)}(p).
	\]
\end{definition}
\symidx{\(\rho_{x_0,t_0}(y,t)\)：平均曲率流单调公式中的后向高斯权}
\symidx{\(\Phi_{x_0,t_0}(t)\)：平均曲率流的高斯加权面积}

观察半径是 \(\sqrt\tau\)：当 \(|x-x_0|\) 远大于这个长度时，高斯权迅速衰减；当时间接近 \(t_0\) 时，同一积分转而关注越来越小的区域。为比较不同尺度，取 \(\lambda>0\)，令
\begin{equation}\label{g3:eq:mcf-parabolic-rescaling}
	F^{(\lambda)}(p,s)
	=\lambda\bigl(F(p,t_0+s/\lambda^2)-x_0\bigr),\qquad s<0.
\end{equation}
度量变为 \(\lambda^2g\)，面积元变为 \(\lambda^m\mu_g\)，迹平均曲率向量变为 \(\lambda^{-1}\mathcal H\)。时间导数也乘 \(\lambda^{-1}\)，故它仍满足平均曲率流，且
\begin{equation}\label{g3:eq:mcf-gaussian-scale-invariance}
	\Phi^{(\lambda)}_{0,0}(s)
	=\Phi_{x_0,t_0}(t_0+s/\lambda^2).
\end{equation}

\subsection{微分和分部积分}
\label{g3:subsec:2803:02}
计算中须区分固定空间点的偏导数与沿移动子流形的导数。以下简记 \(y=F-x_0\)、\(\rho=\rho_{x_0,t_0}(F,t)\)，并以 \(y^\top,y^\perp\) 表示相对于当前切空间的分解。

\begin{lemma}\label{g3:lem:mcf-gaussian-identity}
	固定 \(x_0\in\mathbb R^N\)、\(t_0\in\mathbb R\)，记 \(\rho_{x_0,t_0}(y,t)=[4\pi(t_0-t)]^{-m/2}\exp(-|y-x_0|^2/[4(t_0-t)])\)，\(g(t)=F_t^*g_{\mathrm E}\)、\(\mathcal H=\operatorname{tr}_{g(t)}B\)。在下式中令 \(y=F-x_0\)、\(\tau=t_0-t\)、\(\rho=\rho_{x_0,t_0}(F,t)\)，上标 \(\top,\perp\) 分别表示沿当前浸入的切向与法向投影。
	设 \(F:M^m\times I\to\mathbb R^N\) 为光滑平均曲率流，\(I\subset(-\infty,t_0)\)。则逐点有
	\begin{equation}\label{g3:eq:mcf-gaussian-divergence}
		\frac{\mathrm d}{\mathrm dt}\rho-|\mathcal H|^2\rho
		=\operatorname{div}_{g(t)}\!\left(\frac{\rho y^\top}{2\tau}\right)
		-\rho\left|\mathcal H+\frac{y^\perp}{2\tau}\right|^2.
	\end{equation}
\end{lemma}
\begin{proof}
	在固定环境坐标点求导，得到
	\[
		\partial_t\rho_{x_0,t_0}
		=\left(\frac m{2\tau}-\frac{|y|^2}{4\tau^2}\right)\rho,
		\qquad D\rho_{x_0,t_0}=-\frac{y}{2\tau}\rho.
	\]
	沿 \(F_t=\mathcal H\) 加入空间运动，故
	\[
		\frac{\mathrm d}{\mathrm dt}\rho
		=\left(\frac m{2\tau}-\frac{|y|^2}{4\tau^2}
		-\frac{\langle y,\mathcal H\rangle}{2\tau}\right)\rho.
	\]
	在一点取正交切标架 \(e_1,\dots,e_m\)。因为 \(D_{e_i}y=e_i\)，而微分 \(\langle y^\perp,e_i\rangle=0\) 给
	\(\sum_i\langle D_{e_i}y^\perp,e_i\rangle=-\langle y,\mathcal H\rangle\)，所以
	\[
		\operatorname{div}_M y^\top=m+\langle y,\mathcal H\rangle,
		\qquad \nabla_M\rho=-\frac{\rho y^\top}{2\tau}.
	\]
	乘积法则于是给出
	\[
		\operatorname{div}_M\left(\frac{\rho y^\top}{2\tau}\right)
		=\rho\left(\frac m{2\tau}
		+\frac{\langle y,\mathcal H\rangle}{2\tau}
		-\frac{|y^\top|^2}{4\tau^2}\right).
	\]
	从沿流导数减去此式及 \(|\mathcal H|^2\rho\)，利用 \(\mathcal H\) 法向以及 \(|y|^2-|y^\top|^2=|y^\perp|^2\)，所得三项正好合成所示负平方。
\end{proof}

这里采用通常的正散度，与本书 \(\Delta=-\operatorname{div}\nabla\) 的约定相容。分部积分消去切向项，剩余平方只含法向速度与径向收缩速度的差。

\begin{theorem}[Huisken 单调公式]\label{g3:thm:huisken-monotonicity}
	下式 \(\mathcal H\) 是迹平均曲率向量，\(g(t)=F_t^*g_{\mathrm E}\)，\(\rho_{x_0,t_0}(y,t)=[4\pi(t_0-t)]^{-m/2}e^{-|y-x_0|^2/[4(t_0-t)]}\)，\(\Phi_{x_0,t_0}(t)=\int_M\rho_{x_0,t_0}(F,t)\,\mathrm d\mu_{g(t)}\)。
	设 \(M^m\) 是闭光滑流形，\(m\ge1\)，\(F:M\times[0,T)\to\mathbb R^N\) 是光滑平均曲率流。对任何 \(x_0\in\mathbb R^N\)、\(t_0>0\) 及 \(0\le t<\min(T,t_0)\)，有
	\begin{equation}\label{g3:eq:huisken-monotonicity}
		\frac{\mathrm d}{\mathrm dt}\Phi_{x_0,t_0}(t)
		=-\int_M\rho_{x_0,t_0}(F,t)
		\left|\mathcal H+\frac{(F-x_0)^\perp}{2(t_0-t)}\right|^2\mu_{g(t)}.
	\end{equation}
	特别地，\(\Phi_{x_0,t_0}\) 非增。
\end{theorem}
\begin{proof}
	闭性保证积分可微。由\cref{g3:prop:mcf-area-dissipation}的面积元公式，导数是
	\(\int_M(\mathrm d\rho/\mathrm dt-|\mathcal H|^2\rho)\mu_g\)。代入\cref{g3:lem:mcf-gaussian-identity}，闭流形上的散度积分为零。该计算只用面积密度与散度定理，无须给 \(M\) 选择定向。
\end{proof}

\term[Huisken dandiao gongshi]{Huisken 单调公式}[Huisken monotonicity formula]将积分量的变化写成几何误差的平方。带边界时，式\eqref{g3:eq:mcf-gaussian-divergence}还留下
\[
	\int_{\partial M}\frac{\rho\langle F-x_0,\mathrm dF(\eta)\rangle}{2\tau}\,\mu_{\partial g},
\]
其中 \(\eta\) 是源内的外单位余法向。即使未加权面积的边界通量为零，这个新边界项也未必消失。非紧流则须控制无穷远通量。本节随后使用的非紧圆柱与平面模型，可以直接计算其积分。

\subsection{平方项与等号条件}
\label{g3:subsec:2803:03}
单调公式还指出在什么形状上停止下降。

\begin{proposition}\label{g3:prop:mcf-monotonicity-equality}
	设 \(M^m\) 闭，\(F:M\times[0,T)\to\mathbb R^N\) 为光滑平均曲率流，\(x_0\in\mathbb R^N\)、\(t_0>0\)。以\cref{g3:def:mcf-gaussian-area}的高斯权定义 \(\Phi_{x_0,t_0}\)，并以 \(\mathcal H\) 记迹平均曲率向量。\(\Phi_{x_0,t_0}\) 在非退化时间区间 \([a,b]\subset[0,\min(T,t_0))\) 上恒定，当且仅当在该区间上
	\begin{equation}\label{g3:eq:mcf-shrinking-condition}
		\mathcal H=-\frac{(F-x_0)^\perp}{2(t_0-t)}.
	\end{equation}
	这时经过随时间变化的重参数化，各时间片是某一固定浸入以 \(x_0\) 为中心的同心缩放，缩放比为 \(\sqrt{(t_0-t)/(t_0-a)}\)。
\end{proposition}
\begin{proof}
	积分式\eqref{g3:eq:huisken-monotonicity}表明，恒定性等价于时空中的非负连续被积函数处处为零；因 \(\rho>0\)，即得条件。反向直接代入。
	设 \(\widehat F=(F-x_0)/\sqrt\tau\)。利用该条件，
	\[
		\partial_t\widehat F
		=\frac{\mathcal H}{\sqrt\tau}+\frac{F-x_0}{2\tau^{3/2}}
		=\frac{(F-x_0)^\top}{2\tau^{3/2}}.
	\]
	右端是 \(\widehat F\) 的切向场，可唯一写成 \(\mathrm d\widehat F(Z_t)\)。在闭参数流形上解 \(\dot\psi_t=-Z_t\circ\psi_t\)、\(\psi_a=\operatorname{id}\)，便有 \(\partial_t(\widehat F_t\circ\psi_t)=0\)。恢复尺度即得结论。
\end{proof}

在单个时刻导数为零，只先得到该时间片满足式\eqref{g3:eq:mcf-shrinking-condition}；区间上的几何描述使用区间中的等号。

\begin{definition}\label{g3:def:mcf-self-shrinker}
	光滑浸入 \(X:\Sigma^m\to\mathbb R^N\) 若满足
	\[
		\mathcal H_X+\frac{X^\perp}{2}=0,
	\]
	则称为以原点为中心、尺度归一化的\term[zi shousuo zi]{自收缩子}[self-shrinker]。族 \(\sqrt{-t}\,X\)、\(t<0\)，具有平均曲率法向速度；在能够解相应切向重参数化方程的区间上，它给出纯法向参数化的平均曲率流。
\end{definition}

\begin{proposition}\label{g3:prop:mcf-gaussian-critical}
	设 \(X:\Sigma^m\to\mathbb R^N\) 是闭光滑浸入，\(g=X^*g_{\mathrm E}\) 为诱导度量。它是自收缩子，当且仅当它对高斯面积
	\[
		\mathcal F(X)=(4\pi)^{-m/2}\int_\Sigma e^{-|X|^2/4}\,\mu_g
	\]
	的全部光滑法向变分驻值。
\end{proposition}
\begin{proof}
	对法向初速度 \(V\)，权函数变分为 \(-\tfrac12\langle X,V\rangle e^{-|X|^2/4}\)，面积元变分为 \(-\langle\mathcal H,V\rangle\mu_g\)。因此
	\[
		\delta\mathcal F_X(V)
		=-(4\pi)^{-m/2}\int_\Sigma e^{-|X|^2/4}
		\left\langle\mathcal H+\frac{X^\perp}{2},V\right\rangle\mu_g.
	\]
	任意性给出结论；具体取 \(V=\mathcal H+X^\perp/2\)，驻值就迫使它为零。
\end{proof}

\subsection{自收缩子和奇点放大}
\label{g3:subsec:2803:04}
\begin{example}\label{g3:exm:mcf-shrinker-models}
	经过原点的 \(m\) 维平面有 \(\mathcal H=X^\perp=0\)，是静止自收缩子。对 \(1\le k\le m\)，圆柱
	\[
		S^k(\sqrt{2k})\times\mathbb R^{m-k}\subset\mathbb R^{m+1}
	\]
	也满足自收缩方程：写 \(X=(z,w)\)，则 \(X^\perp=(z,0)\)，球面部分给 \(\mathcal H=(-kz/|z|^2,0)=-(z,0)/2\)。相应流为 \(S^k(\sqrt{-2kt})\times\mathbb R^{m-k}\)。\(k=m\) 时得到闭球面，\(k<m\) 时是非紧模型，轴向高斯积分仍有限。上述球面及非平坦圆柱模型的曲率平方为 \(1/[2(-t)]\)。
\end{example}

设闭光滑流的存在区间为 \([0,T)\)，\(T<\infty\)。由单调性与非负性，对每个 \(x_0\) 都存在有限极限
\begin{equation}\label{g3:eq:mcf-gaussian-density}
	\Theta(x_0,T)\coloneqq\lim_{t\uparrow T}\Phi_{x_0,T}(t).
\end{equation}
它称为该时空点的\term[gaosi midu]{高斯密度}[Gaussian density]。极限存在只需要标量单调性；要从放大的流得到光滑几何极限，还需控制导数及局部面积。下面先把已有光滑极限必须满足的方程证出。

\begin{proposition}\label{g3:prop:mcf-smooth-blowup-shrinker}
	设 \(F:M^m\times[0,T)\to\mathbb R^N\) 是闭光滑平均曲率流，\(T<\infty\)。固定 \(x_0\)，取 \(\lambda_j\to\infty\)，按式\eqref{g3:eq:mcf-parabolic-rescaling}以 \(t_0=T\) 放大。若这些流在某开时空区域中经过局部重参数化光滑收敛，且局部片的重数为有限正整数，则每个光滑极限片在 \(s<0\) 上满足
	\[
		\mathcal H_\infty+\frac{X^\perp}{2(-s)}=0.
	\]
\end{proposition}
\begin{proof}
	固定 \(a<b<0\)，取 \(j\) 充分大使相应原时间为正。由尺度不变性与单调公式，
	\begin{align*}
		&\int_a^b\int_M\rho_{0,0}(F^{(\lambda_j)},s)
		\left|\mathcal H_j+\frac{(F^{(\lambda_j)})^\perp}{2(-s)}\right|^2
		\mu_{g_j(s)}\,\mathrm ds\\
		&\qquad=\Phi_{x_0,T}(T+a/\lambda_j^2)
		-\Phi_{x_0,T}(T+b/\lambda_j^2)\longrightarrow0.
	\end{align*}
	右边两个值都趋于 \(\Theta(x_0,T)\)。任意相对紧的光滑极限参数片、紧时间子区间上的非负积分不超过左边。光滑收敛允许在该局部积分中取极限，有限正重数保留其零值。高斯权严格为正，故连续平方项在每个这样的片上为零，任意性给出结论。此处没有交换整个非紧极限上的积分与极限。
\end{proof}

\begin{figure}[htbp]
	\centering
	% !TeX root = ../Book3.tex
\begin{tikzpicture}[x=1cm,y=1cm,>=stealth]
	\draw[->] (0,0)--(3.0,0) node[right] {空间};
	\draw[->] (1.4,-1.9)--(1.4,.6) node[above] {时间};
	\fill[blue!10] (1.0,-.6) rectangle (1.8,0);
	\draw[blue!60!black] (1.0,-.6) rectangle (1.8,0);
	\fill (1.4,0) circle (1.6pt);
	\node[above right] at (1.4,0) {\((x_0,T)\)};
	\node[below] at (1.4,-.65) {\(\lambda^{-1},\ \lambda^{-2}\)};
	\draw[->,thick] (3.2,-.6)--(4.6,-.6) node[midway,above] {放大};
	\draw[->] (4.9,0)--(8.0,0) node[right] {空间};
	\draw[->] (6.3,-1.9)--(6.3,.6) node[above] {时间};
	\fill[blue!10] (5.3,-1.4) rectangle (7.3,0);
	\draw[blue!60!black] (5.3,-1.4) rectangle (7.3,0);
	\fill (6.3,0) circle (1.6pt);
	\node[above right] at (6.3,0) {\((0,0)\)};
	\node[below] at (6.3,-1.45) {固定空间与时间尺度};
\end{tikzpicture}

	\caption{以 \((x_0,T)\) 为中心的抛物放大。空间长度量级 \(\lambda^{-1}\)、时间长度量级 \(\lambda^{-2}\) 的观察区域变为固定尺度；图示只记录尺度关系。}
	\label{g3:fig:mcf-parabolic-rescaling}
\end{figure}

若有限奇异时刻满足 \(\sup_M|B(t)|^2\le C/(T-t)\)，称其满足\term[I xing qidian]{I 型奇点}[type I singularity]的曲率增长界；不满足此界的有限时奇点称为\term[II xing qidian]{II 型奇点}[type II singularity]。取 \(\lambda_j=(T-t_j)^{-1/2}\)，I 型界使放大后 \(s=-1\) 时间片的曲率一致有界。\cref{g3:fig:mcf-parabolic-rescaling}说明时间平方尺度的来源。光滑极限的存在需要进一步紧性论证；一旦有这种极限，前一命题确定其自收缩方程。球面与圆柱展示 I 型尺度，弱极限的重数及奇点后的延续则由下一节分别讨论。
% 纲要条目排为编号小节标题，并进入目录与书签；说明段仍保留于纲要注释中。

% !TeX root = ../../Book3.tex
% 纲要标题：### 28.4 弱演化的不同对象
% 纲要位置：几何学计划纲要.md，第 3306 行。
% 正文状态：初稿；尚未正式审核。

% 纲要内容开始
% 
% * 水平集方程及粘性解
% * Brakke 不等式及测度族
% * 质量丢失、重数和肥化
% * 存在、唯一性及正则性的不同结论范围
% 
% ---
% 
% 纲要内容结束

\section{弱演化的不同对象}
\label{g3:sec:2804}

圆球面缩到一点时，浸入及其诱导度量同时退化。继续研究演化，需要先决定保留哪些信息：水平集方法保留区域及其边界的位置，变分流形方法还保留面积和重数。这两种选择产生不同的解空间，也使“唯一”具有不同含义。本节分别给出方程与定义，证明它们在光滑区域中的相容性，再陈述所用的整体存在结果。

\subsection{水平集方程及粘性解}
\label{g3:subsec:2804:01}
设 \(u=u(x,t)\) 是 \(\mathbb R^n\times(0,T)\) 上的光滑函数，\(n\ge2\)，并在所考察的零水平集上满足 \(\mathrm Du\ne0\)。选择法向 \(\nu=\mathrm Du/|\mathrm Du|\)，则该水平集的迹平均曲率向量为
\[
\mathcal H=-\operatorname{div}\left(\frac{\mathrm Du}{|\mathrm Du|}\right)\nu.
\]
如果 \(x(t)\) 沿水平集以速度 \(\mathcal H\) 移动，对 \(u(x(t),t)=0\) 求导得到
\begin{equation}\label{g3:eq:mcf-level-set}
	u_t=|\mathrm Du|\operatorname{div}\left(\frac{\mathrm Du}{|\mathrm Du|}\right)
	=\left(\delta_{ij}-\frac{u_i u_j}{|\mathrm Du|^2}\right)u_{ij}.
\end{equation}
这里使用环境坐标的一阶、二阶导数；右端是水平集切方向上的 Hessian 迹。它在法向方向没有二阶扩散，且在 \(\mathrm Du=0\) 时没有单一的经典取值。因而\cref{g3:thm:quasilinear-parabolic-short-time}的强抛物假设在这里不成立。

\ProseParagraphBreak
粘性解用从上、从下接触的光滑函数测试方程。先写出零梯度处应采用的两个极限。对 \(p\ne0\) 和实对称 \(n\times n\) 矩阵 \(X\)，记
\[
\mathcal F(p,X)=\operatorname{tr}X-\frac{\langle Xp,p\rangle}{|p|^2}.
\]
令 \(\mathcal F^*\)、\(\mathcal F_*\) 分别为当 \((q,Y)\to(p,X)\)、\(q\ne0\) 时的上极限与下极限。在 \(p\ne0\) 处它们都等于 \(\mathcal F\)；在 \(p=0\) 处，Rayleigh 商的极值给出
\begin{equation}\label{g3:eq:mcf-level-set-envelopes}
	\mathcal F^*(0,X)=\operatorname{tr}X-\lambda_{\min}(X),
	\qquad
	\mathcal F_*(0,X)=\operatorname{tr}X-\lambda_{\max}(X).
\end{equation}
事实上，任意非零 \(q\) 的方向都在单位球面上；取趋于最小或最大特征向量的方向，就分别达到上述两端。

\begin{definition}\label{g3:def:mcf-viscosity-solution}
	设 \(Q\subset\mathbb R^n\times(0,\infty)\) 开，\(u:Q\to\mathbb R\) 连续。如果对每个 \(\varphi\in C^{2,1}(Q)\) 及每个使 \(u-\varphi\) 取得局部最大值的 \((x_0,t_0)\in Q\)，都有
	\[
	\varphi_t(x_0,t_0)
	\le\mathcal F^*\bigl(\mathrm D\varphi(x_0,t_0),\mathrm D^2\varphi(x_0,t_0)\bigr),
	\]
	则称 \(u\) 为式\eqref{g3:eq:mcf-level-set}的\term[nianxing cixie]{粘性次解}[viscosity subsolution]。若在局部最小值点满足反向不等式，且右端使用 \(\mathcal F_*\)，就称为\term[nianxing chaoxie]{粘性超解}[viscosity supersolution]。同时满足两种条件的连续函数称为\term[nianxing jie]{粘性解}[viscosity solution]。初值 \(u_0\) 通过 \(u(\cdot,t)\to u_0\) 的连续延拓施加。
\end{definition}

上述定义按 \(u_t-\mathcal F=0\) 的符号写出；次解使用右端的上包络，超解使用下包络。它是文献\cite[定义2.1，第753--754页]{ChenGigaGoto1991MeanCurvature}对奇异算子的定义在本方程中的具体形式。条件测试的是所有接触函数，并非要求 \(u\) 本身处处可微。

\begin{proposition}\label{g3:prop:mcf-level-set-smooth-consistency}
	设 \(u\in C^{2,1}(Q)\)，其中 \(Q\subset\mathbb R^n\times(0,T)\) 开，且 \(\mathrm Du\ne0\) 在 \(Q\) 上成立。则 \(u\) 是式\eqref{g3:eq:mcf-level-set}的粘性解，当且仅当它是经典解；此时每个非空水平集的法向速度都是其迹平均曲率向量。
\end{proposition}
\begin{proof}
	若 \(u\) 是粘性解，直接取测试函数 \(\varphi=u\)，次解与超解条件给出等式。反之，在 \(u-\varphi\) 的局部最大值点，\(u_t=\varphi_t\)、\(\mathrm Du=\mathrm D\varphi\)，且 \(\mathrm D^2u\le\mathrm D^2\varphi\)。矩阵 \(I-p\otimes p/|p|^2\) 半正定，所以 \(\mathcal F(p,\mathrm D^2u)\le\mathcal F(p,\mathrm D^2\varphi)\)，得到次解条件；局部最小值处同理。

	对水平集上一条运动曲线，有 \(u_t+\mathrm Du\cdot x'=0\)，故其法向速度是 \(-u_t\nu/|\mathrm Du|\)。代入方程恰为 \(\mathcal H\)。也可直接求解 \(x'=\mathcal H(x,t)\) 的局部常微分方程；沿解 \(u\) 保持常值，从而得到以纯法向速度移动的局部参数。
\end{proof}

\begin{proposition}\label{g3:prop:mcf-level-set-relabel}
	设 \(u\) 是开时空区域 \(Q\) 上的连续粘性解，\(\chi\in C^2(\mathbb R)\) 且 \(\chi'>0\)。则 \(v=\chi(u)\) 也是粘性解。
\end{proposition}
\begin{proof}
	设 \(v-\varphi\) 在 \((x_0,t_0)\) 取得局部最大值，平移测试函数使两者在该点相等。\(\chi\) 在该值附近可逆，\(\psi=\chi^{-1}\circ\varphi\) 从上接触 \(u\)。记该点 \(a=\chi'(u)>0\)、\(b=\chi''(u)\)，则
	\[
	\varphi_t=a\psi_t,\qquad
	\mathrm D\varphi=a\mathrm D\psi,\qquad
	\mathrm D^2\varphi=a\mathrm D^2\psi+b\,\mathrm D\psi\otimes\mathrm D\psi.
	\]
	对 \(p\ne0\)，\(\mathcal F(ap,aX+bp\otimes p)=a\mathcal F(p,X)\)，因为切向投影消去 \(p\otimes p\)。对 \(p=0\)，由式\eqref{g3:eq:mcf-level-set-envelopes}，两个包络也满足相应的正齐次恒等式。因此 \(u\) 的次解条件给出 \(v\) 的次解条件。从下接触时相同计算给出超解条件。
\end{proof}

若 \(\chi(0)=0\)，这一变换保持零水平集以及两侧符号；因此描述函数的数值标记具有自由度。任意连续定义函数之间的独立性还需要比较原理，将在本节末准确引用。

\subsection{Brakke 不等式及测度族}
\label{g3:subsec:2804:02}
接着保留面积信息。按\cref{g3:def:rectifiable-varifold}，整变分流形可写成 \(V=\mathbf v(S,\theta)\)，权测度为 \(\mu=\theta\mathcal H^m\llcorner S\)，其中 \(\theta\) 几乎处处取正整数。若第一变分没有奇异部分，则\eqref{g3:eq:varifold-mean-curvature-decomposition}成为
\begin{equation}\label{g3:eq:brakke-mean-curvature}
	\delta V(X)=-\int\mathcal H_V\cdot X\,\mathrm d\mu,
	\qquad X\in C_c^1(\mathbb R^n;\mathbb R^n).
\end{equation}
这里的 \(\mathcal H_V\) 就是前章记作 \(H_V\) 的广义平均曲率；对光滑 \(m\) 维子流形，它等于 \(mH=\mathcal H\)，这一换算已由\cref{g3:prop:smooth-varifold-boundary-variation}证明。

\begin{proposition}\label{g3:prop:mcf-weighted-area-transport}
	设 \(F:M^m\times[a,b]\to\mathbb R^n\) 是闭光滑流形上的光滑平均曲率流，速度为 \(F_t=\mathcal H\)。对 \(\varphi\in C^{2,1}(\mathbb R^n\times[a,b])\)，有
	\begin{equation}\label{g3:eq:mcf-weighted-area-transport}
		\frac{\mathrm d}{\mathrm dt}\int_M\varphi(F,t)\,\mu_{g(t)}
		=\int_M\bigl(\varphi_t+\mathrm D\varphi\cdot\mathcal H
		-\varphi|\mathcal H|^2\bigr)(F,t)\,\mu_{g(t)}.
	\end{equation}
\end{proposition}
\begin{proof}
	沿运动点的函数导数为 \(\varphi_t+\mathrm D\varphi\cdot F_t\)。面积元导数为 \(-|\mathcal H|^2\mu_{g(t)}\)。将两项相加，并利用 \(M\) 紧以交换求导与积分，即得结论。
\end{proof}

弱极限中，面积可能丢失，而曲率平方积分具有下半连续性。将上式等号改成只允许额外下降的不等号，便得到 Brakke 演化。\footnote{Kenneth A. Brakke，美国数学家，研究几何测度论与曲面的平均曲率演化，并开发了 Surface Evolver 曲面计算软件；人物与软件资料见其\href{https://kenbrakke.com/}{个人主页}。}下面采用带时间测试函数的积分形式，固定端点测度的取值，从而把跳跃下降也纳入定义。

\begin{definition}\label{g3:def:integral-brakke-flow}
	设 \(1\le m<n\)、\(I=[0,T)\)，其中 \(T\le\infty\)。一族 \(\mathbb R^n\) 上的非负 Radon 测度 \(\{\mu_t\}_{t\in I}\) 称为无边界的 \(m\) 维\term[Brakke liu]{整 Brakke 流}[integral Brakke flow]，若满足下列条件：
	\begin{enumerate}
		\item 对每个紧集 \(K\subset\mathbb R^n\) 和 \(S<T\)，有 \(\sup_{0\le t\le S}\mu_t(K)<\infty\)；对每个 \(f\in C_c(\mathbb R^n)\)，\(t\mapsto\mu_t(f)\) 可测。
		\item 对几乎每个 \(t\)，\(\mu_t\) 是整变分流形 \(V_t\) 的权测度，\(V_t\) 的第一变分满足式\eqref{g3:eq:brakke-mean-curvature}。其广义平均曲率 \(\mathcal H_t\) 垂直于近似切平面，并可选成关于 \(\mathrm d\mu_t\,\mathrm dt\) 可测的向量场，满足
		\[
		\int_0^S\int_K|\mathcal H_t|^2\,\mathrm d\mu_t\,\mathrm dt<\infty
		\quad(K\text{ 紧},\ S<T).
		\]
		\item 对每个 \(0\le a<b<T\) 和每个空间支集一致紧的非负 \(\varphi\in C^{2,1}(\mathbb R^n\times[a,b])\)，都有
		\begin{equation}\label{g3:eq:brakke-integrated-inequality}
			\mu_b(\varphi(\cdot,b))-\mu_a(\varphi(\cdot,a))
			\le\int_a^b\int_{\mathbb R^n}
			\bigl(\varphi_t+\mathrm D\varphi\cdot\mathcal H_t
			-\varphi|\mathcal H_t|^2\bigr)\,\mathrm d\mu_t\,\mathrm dt.
		\end{equation}
	\end{enumerate}
	称式\eqref{g3:eq:brakke-integrated-inequality}为\term[Brakke budengshi]{Brakke 不等式}[Brakke inequality]。说该流连续地取初值 \(\mu_0\)，另指 \(\mu_t\rightharpoonup\mu_0\) 当 \(t\downarrow0\)，即所有紧支连续函数上的积分收敛。
\end{definition}

这里 \(C^{2,1}\) 指空间二阶、时间一阶连续可微。局部质量界及平方曲率的时空可积性使右端各项可积；例如 \(\mathrm D\varphi\cdot\mathcal H_t\) 可由 Cauchy--Schwarz 控制。在异常的零测时间集上无需定义曲率，端点测度却仍有明确取值。不等式要求每对端点都成立，因而在某个单独时刻任意增加质量通常会破坏定义。原始的上导数形式见\cite[第3.2--3.3节，第20--22页]{Brakke1978Motion}，其构造所满足的积分不等式见该书第4.29节，式(1)。

\begin{corollary}\label{g3:cor:smooth-flow-is-brakke}
	闭光滑嵌入子流形的光滑平均曲率流，其面积测度构成整 Brakke 流，并在式\eqref{g3:eq:brakke-integrated-inequality}中取等号。若一整 Brakke 流在 \([a,b]\) 的全部支集包含于同一紧集，则
	\begin{equation}\label{g3:eq:brakke-total-dissipation}
		\mu_b(\mathbb R^n)+\int_a^b\int|\mathcal H_t|^2\,\mathrm d\mu_t\,\mathrm dt
		\le\mu_a(\mathbb R^n).
	\end{equation}
\end{corollary}
\begin{proof}
	光滑情形的第一变分和曲率均由前述公式给出，其余条件由紧性成立。对式\eqref{g3:eq:mcf-weighted-area-transport}在时间上积分即得等号。第二个结论中，取一个在所给紧集的邻域等于一的非负紧支空间函数作为 \(\varphi\)，则其时间导数及支集上的空间梯度都为零，故得到所示不等式。
\end{proof}

\subsection{质量丢失、重数和肥化}
\label{g3:subsec:2804:03}
Brakke 不等式保留光滑流的耗散方向，但允许更多的质量下降。下面的例子在曲率仍然有界时就能显示这一自由度。

\begin{example}\label{g3:exm:brakke-premature-extinction}
	取半径 \(R_0>0\) 的 \(m\) 维圆球面，令
	\[
	R(t)=\sqrt{R_0^2-2mt},\qquad T_*=\frac{R_0^2}{2m}.
	\]
	固定 \(0<c<T_*\)，定义
	\[
	\mu_t^{(c)}=
	\begin{cases}
		\mathcal H^m\llcorner S^m_{R(t)},&0\le t<c,\\
		0,&t\ge c.
	\end{cases}
	\]
	这是一个连续地取同一初值的整 Brakke 流。确实，\(t<c\) 时它满足光滑等式；\(t\ge c\) 时各项都为零。若 \(a<c\le b\)，把光滑等式积分到 \(c\) 的左侧，右端等于
	\[
	\int_{S^m_{R(c)}}\varphi(x,c)\,\mathrm d\mathcal H^m(x)
	-\mu_a^{(c)}(\varphi(\cdot,a)),
	\]
	它不小于左端 \(-\mu_a^{(c)}(\varphi(\cdot,a))\)。因为 \(\varphi\ge0\)，跳跃恰好沿允许的方向发生。曲率在 \([0,c]\) 有界，故平方曲率可积条件也满足。不同的 \(c\) 给出同一初始球面的不同弱流；这一差别来自定义允许的\term[zhiliang diushi]{质量丢失}[mass drop]。
\end{example}

\begin{example}\label{g3:exm:brakke-multiplicity}
	设 \(\mu_t\) 是整 Brakke 流，\(q\) 是固定正整数。则 \(q\mu_t\) 也是整 Brakke 流，并有相同的广义平均曲率。因为第一变分和权测度同时乘以 \(q\)，式\eqref{g3:eq:brakke-mean-curvature}中的曲率向量不变；不等式两端也同时乘以 \(q\)。两者有相同支集，却有不同面积和重数。例如重数一与重数二的收缩球面具有同样的几何位置。这是在比较不同初始测度，不能据此说固定初始测度已经不唯一；上一例才提供了后一种现象。
\end{example}

水平集方法中的另一种退化是零水平集变厚。设连续解 \(u\) 的初始零水平集没有内点。如果某个 \(t>0\) 满足
\[
\operatorname{int}\{x:u(x,t)=0\}\ne\varnothing,
\]
就称该零水平集在时刻 \(t\) 发生\term[feihua]{肥化}[fattening]。此时零水平集含有开区域，不能作为通常的余维一界面理解。若记
\[
\Omega_t^- =\{u(\cdot,t)<0\},\qquad K_t=\{u(\cdot,t)\le0\},
\]
则 \(K_t\setminus\Omega_t^-=\{u=0\}\) 可以含有开集；由严格负相和非正相提取界面时，需要说明采用哪种边界。

\ProseParagraphBreak
肥化描述的是一个确定的函数解的水平集几何。即使函数解唯一，集合 \(\{u=0\}\) 仍可能变厚；这两个断言并不矛盾。相反，若解在零水平集附近光滑且梯度不为零，隐函数定理保证该水平集局部是超曲面，从而没有内点。质量丢失、整数重数与肥化分别涉及测度随时间的跳跃、面积的计数和零集的空间结构，判断它们所需的信息各不相同。

\subsection{存在、唯一性及正则性的不同结论范围}
\label{g3:subsec:2804:04}
水平集方程的比较原理必须处理零梯度处的不连续算子；Brakke 流的构造则要同时保持测度紧性、整数重数和平方曲率控制。本节外引这两项整体理论，给出便于应用的具体版本。

\begin{theorem}\label{g3:thm:level-set-global-external}
	设 \(n\ge2\)，\(u_0:\mathbb R^n\to\mathbb R\) 连续，且在某个紧集外等于常数 \(c\)。则式\eqref{g3:eq:mcf-level-set}存在唯一连续粘性解 \(u\) 于 \(\mathbb R^n\times[0,\infty)\)，取初值 \(u_0\)，并具有如下性质：对每个有限 \(S>0\)，有 \(R_S<\infty\) 使 \(u(x,t)=c\) 当 \(|x|>R_S\)、\(0\le t\le S\)。唯一性在这一解类中成立。

	若 \(u_0\le v_0\)，且两者分别在紧集外为常数，则相应解满足 \(u\le v\)。若 \(c>0\)，则集合对
	\[
	\bigl(\{u(\cdot,t)=0\},\ \{u(\cdot,t)<0\}\bigr)
	\]
	只由初始的这一集合对决定，与所用的连续定义函数无关。
\end{theorem}

这一结果取自\cite[定理6.7、6.8及注6.9，第780--781页；定理7.1，第782--783页]{ChenGigaGoto1991MeanCurvature}，证明分别使用奇异算子的比较原理、上下解构造以及保持水平集的函数变换。原文把正相作为有界区域；对本书的负相约定，将 \(u\) 换成 \(-u\) 即可。方程在此变换下不变，因为 \(\mathcal F(-p,-X)=-\mathcal F(p,X)\)，零梯度处的两个包络随之互换。初值在紧集外为常数还保证它有界且一致连续，本定理并未把结论的解类扩大到任意无穷远增长。

\begin{corollary}\label{g3:cor:level-set-contraction}
	在\cref{g3:thm:level-set-global-external}的解类中，对任意两个初值 \(u_0,v_0\) 及任意 \(t\ge0\)，有
	\[
	\|u(\cdot,t)-v(\cdot,t)\|_{L^\infty(\mathbb R^n)}
	\le\|u_0-v_0\|_{L^\infty(\mathbb R^n)}.
	\]
\end{corollary}
\begin{proof}
	令右端为 \(d\)。方程不含函数值本身，故 \(v+d\) 和 \(v-d\) 仍在指定类中求解。初值满足 \(v_0-d\le u_0\le v_0+d\)，两次应用比较原理即可。
\end{proof}

测度演化有如下存在结论，其证明属于变分流形的逼近理论，见\cite[第4.26节，第106--108页；第4.29节，第116--120页]{Brakke1978Motion}。这里只取光滑紧初始超曲面的版本。

\begin{theorem}\label{g3:thm:brakke-existence-external}
	设 \(\Sigma_0^m\subset\mathbb R^{m+1}\) 是非空、闭、光滑嵌入超曲面，\(m\ge1\)。存在定义于全部 \(t\ge0\) 的 \(m\) 维整 Brakke 流 \(\{\mu_t\}\)，使
	\[
	\mu_0=\mathcal H^m\llcorner\Sigma_0,
	\qquad \mu_t\rightharpoonup\mu_0\quad(t\downarrow0).
	\]
\end{theorem}

原构造先平滑第一变分，再以小时间步推进，并在每步加入不增加所测试面积的局部变形；局部质量界给出弱极限，曲率控制和整性紧性使几乎每个时间片恢复整变分流形。所述定理选出了一个连续地取给定初值的整体弱演化。若要比较该构造与经典解在整个光滑区间内的一致性，还需要所选逼近的相容性质；仅满足 Brakke 不等式的任意流可以像\cref{g3:exm:brakke-premature-extinction}那样提前丢失质量。

\ProseParagraphBreak
两个存在结论都允许越过经典浸入的奇点，但输出不同。水平集定理确定连续函数解及相应的集合演化，描述中没有整数重数；Brakke 定理输出带面积与重数的测度族，定义允许质量下降，因而没有一般唯一性结论。在二者的光滑重合区域，\cref{g3:prop:mcf-level-set-smooth-consistency,g3:cor:smooth-flow-is-brakke}给出了同一个法向速度与面积演化。

\ProseParagraphBreak
正则性则要求从弱对象重新得到光滑参数、切平面和曲率控制。单调公式中的非负平方项提供尺度间的耗散信息，零耗散给出自相似约束；要将小耗散或密度信息提升为光滑性，仍须另有相应的正则性定理及其密度、重数或非肥化假设。本章的整体外引结论是上述水平集比较与存在、以及指定初值的整 Brakke 流存在；有关所有奇点的光滑性或奇点后经典解的唯一延续，均不包含在这些结论之中。
% 纲要条目排为编号小节标题，并进入目录与书签；说明段仍保留于纲要注释中。


\begin{chaptersummary}
	平均曲率流的速度取 \(\mathcal H=\operatorname{tr}_gB=mH\)。参数方程的主符号在切向方向退化；把邻近浸入写成初始法丛上的图，法向投影与竖直微分的同构使最高阶部分成为 \(g_u^{ij}\nabla_i\nabla_j\)。第26章的短时理论和去除切向速度的常微分方程共同给出闭浸入的短时存在与参数化唯一性。

	\ProseParagraphBreak
	面积下降控制曲率迹的时空平方积分。超曲面的第二基本形式还满足含立方反应项的抛物方程，标量比较与二次型的首次零方向论证分别保持平均凸性和凸性。避免相交使用距离极小处两个切空间相同这一余维一事实；正交平面内的两个缩圆说明高余维时结论可以失效。

	\ProseParagraphBreak
	高斯权采用子流形维数，因而加权面积在抛物缩放下不变。微分后以散度项消去切向贡献，得到 Huisken 单调公式。等号对应自收缩方程；光滑放大极限若存在，其局部耗散为零，便满足同一方程。这个结论把极限的识别与极限的存在分开。

	\ProseParagraphBreak
	水平集方法以函数的零集描述演化，粘性定义在零梯度处使用算子的半连续包络。Brakke 方法以带重数的面积测度和广义平均曲率描述演化，保留积分形式的面积下降不等式。粘性函数的唯一性、零集是否肥化、Brakke 流是否丢失质量以及弱流何处光滑，是不同的问题；缩球与提前消失的例子使这些区别可以直接核验。
\end{chaptersummary}

\BookExercises

\begin{exercise}\label{g3:ex:mcf-product-spheres}
	设 \(p,q\ge1\)，在 \(\mathbb R^{p+1}\times\mathbb R^{q+1}\) 中考虑 \(F(x,y,t)=(r(t)x,R(t)y)\)，其中 \(x\in S^p\)、\(y\in S^q\)，初始半径均为正。求使其成为平均曲率流的半径方程与光滑存在时间，判断何时它是同心自相似收缩，并验证面积耗散。
\end{exercise}
\begin{solution}
	乘积度量为 \(r^2g_{S^p}\oplus R^2g_{S^q}\)，两组第二基本形式分别取值于两个正交法向，所以
	\[
		\mathcal H=(-p x/r,-q y/R),\qquad
		r^2=r_0^2-2pt,\quad R^2=R_0^2-2qt.
	\]
	光滑浸入存在至 \(T=\min\{r_0^2/(2p),R_0^2/(2q)\}\)。两半径的比保持常数当且仅当 \(r_0^2/p=R_0^2/q\)，此时同时缩为一点。否则先缩掉一个球因子，另一因子尚有正半径；在此时参数微分已经退化。面积为 \(\sigma_p\sigma_q r^pR^q\)，其中 \(\sigma_k=\operatorname{Vol}(S^k(1))\)，其对数导数为 \(-p^2/r^2-q^2/R^2=-|\mathcal H|^2\)，正是耗散公式。
\end{solution}

\begin{exercise}\label{g3:ex:mcf-symmetry}
	设 \(F_0:M\to(N,\bar g)\) 是闭光滑浸入，\(I:N\to N\) 为等距微分同胚，\(\psi:M\to M\) 为微分同胚，且 \(I\circ F_0=F_0\circ\psi\)。证明其平均曲率流在光滑存在期间满足 \(I\circ F_t=F_t\circ\psi\)。
\end{exercise}
\begin{solution}
	等距映射保持环境联络、诱导度量与法向投影，故 \(\mathcal H_{I\circ F}=\mathrm dI(\mathcal H_F)\)；时间不变的重参数化给 \(\mathcal H_{F\circ\psi}=\mathcal H_F\circ\psi\)。因此 \(I\circ F_t\) 与 \(F_t\circ\psi\) 都是纯法向平均曲率流，并有相同初始映射。\Cref{g3:thm:mcf-short-time}的唯一性使它们在共同存在期间相等。
\end{solution}

\begin{exercise}\label{g3:ex:mcf-periodic-convex-hull}
	设 \(u:\mathbb T^m\times[0,T)\to\mathbb R^q\) 为光滑周期图方程 \(u_t=g^{ij}(\mathrm Du)u_{ij}\) 的解，其中 \(g_{ij}=\delta_{ij}+\sum_\alpha u_i^\alpha u_j^\alpha\)。证明 \(u(\mathbb T^m,t)\) 始终包含在初始像的闭凸包中。
\end{exercise}
\begin{solution}
	任取 \(a\in\mathbb R^q\)，函数 \(v=a\cdot u\) 满足标量方程 \(v_t=g^{ij}v_{ij}\)。在任意紧时间区间，\(\mathrm Du\) 有界，所以主系数一致正定。最大值原理给 \(a\cdot u(x,t)\le\max_y a\cdot u(y,0)\)。这些不等式对每个 \(a\) 成立，其半空间交就是初始像的闭凸包。为核对最后一点，若一点 \(z\) 在紧凸包外，取距它最近的凸包内点 \(z_*\)，则极小性给 \((z-z_*)\cdot(w-z_*)\le0\) 对所有凸包内 \(w\) 成立；方向 \(a=z-z_*\) 因而将 \(z\) 与该凸包严格分离。
\end{solution}

\begin{exercise}\label{g3:ex:mcf-enclosed-volume}
	设闭嵌入超曲面 \(F_t:M^m\to\mathbb R^{m+1}\) 在一个光滑时间区间内为有界区域 \(\Omega_t\) 的边界，取外向单位法向 \(\nu\)，以 \(s=\langle\mathcal H,\nu\rangle\) 记迹平均曲率标量，令 \(\theta=-s\)，流速为 \(-\theta\nu\)。证明
	\[
		\frac{\mathrm d}{\mathrm dt}\operatorname{Vol}(\Omega_t)
		=-\int_{\partial\Omega_t}\theta\,\mu_g.
	\]
	对圆球核对此式，并说明面积与围成体积的耗散为何是两个不同的积分。记 \(\sigma_m=\operatorname{Vol}(S^m(1))\)。
\end{exercise}
\begin{solution}
	将边界速度光滑延拓为邻域内的环境向量场 \(V\)，并取紧支延拓。其局部流给出区域变动的一阶表示，故体积导数为 \(\int_{\Omega_t}\operatorname{div}V\)。散度定理把它化为 \(\int_{\partial\Omega_t}\langle V,\nu\rangle=-\int\theta\)。这个结果只取决于边界法向速度，故不依赖延拓。圆球体积为 \(\sigma_m R^{m+1}/(m+1)\)，导数为 \(\sigma_mR^m(-m/R)\)，与 \(\theta=m/R\) 相符。面积导数是 \(-\int\theta^2\)，始终非正；体积导数是 \(-\int\theta\)，需要曲率符号才能判断其正负。
\end{solution}

\begin{exercise}\label{g3:ex:mcf-umbilic-defect}
	在闭双侧 Euclid 超曲面的光滑平均曲率流上，记 \(M^m\) 为源，\(A\) 为相对于所选单位法向的形算子，\(s=\operatorname{tr}A\)，\(\Delta=\nabla^*\nabla\) 为当前非负粗 Laplace，令 \(D=|A|^2-s^2/m\)。证明 \(D\ge0\)，并推导
	\[
		(\partial_t+\Delta)D
		=-2\left(|\nabla A|^2-\frac{|\nabla s|^2}{m}\right)+2|A|^2D.
	\]
	由此证明初始全脐性在光滑存在期间保持。
\end{exercise}
\begin{solution}
	在正交主方向上，\(D=\sum_i(\kappa_i-s/m)^2\ge0\)，等号恰为 \(A=(s/m)\operatorname{Id}\)。由 \((\partial_t+\Delta)s=|A|^2s\) 与非负 Laplace 的乘积公式，
	\[
		(\partial_t+\Delta)s^2=2|A|^2s^2-2|\nabla s|^2.
	\]
	结合 \(|A|^2\) 的演化得到所求式。对每个正交方向 \(e_j\)，Cauchy 不等式给 \(|\nabla_j s|^2=|\operatorname{tr}\nabla_jA|^2\le m|\nabla_jA|^2\)，求和后有括号非负。因此 \((\partial_t+\Delta)D\le2|A|^2D\)。初值为零时，在任意紧时间区间应用标量最大值原理得 \(D\le0\)，与 \(D\ge0\) 结合即得全脐性保持。
\end{solution}

\begin{exercise}\label{g3:ex:mcf-plane-gaussian}
	设 \(P\subset\mathbb R^N\) 是一个静止的 \(m\) 维仿射平面，\(a=\operatorname{dist}(x_0,P)\)。固定 \(x_0\in\mathbb R^N\)、\(t<t_0\)，以 \(\rho_{x_0,t_0}(x,t)=[4\pi(t_0-t)]^{-m/2}e^{-|x-x_0|^2/[4(t_0-t)]}\) 记后向高斯权。计算 \(\int_P\rho_{x_0,t_0}(x,t)\,\mathrm d\mathcal H^m(x)\)，并直接验证单调公式及等号条件。
\end{exercise}
\begin{solution}
	把 \(x-x_0\) 分成沿平面的向量 \(z\) 与长度为 \(a\) 的固定法向量，Gaussian 积分给
	\[
		\Phi(t)=e^{-a^2/(4\tau)},\qquad
		\Phi'(t)=-\frac{a^2}{4\tau^2}e^{-a^2/(4\tau)}.
	\]
	平面上 \(\mathcal H=0\)、\(|(x-x_0)^\perp|=a\)，所以右式正是单调公式的负平方积分。权恒定当且仅当 \(a=0\)，即中心在平面上。此非紧算例通过直接积分成立，无需预设无穷远分部积分。
\end{solution}

\begin{exercise}\label{g3:ex:mcf-cylinder-density}
	对 \(1\le k\le m\)，计算缩圆柱 \(S^k(\sqrt{-2kt})\times\mathbb R^{m-k}\)、\(t<0\)，以 \((0,0)\) 为中心的高斯加权面积。说明它为何与 \(t\) 及轴向维数 \(m-k\) 无关。记 \(\sigma_k=\operatorname{Vol}(S^k(1))\)。
\end{exercise}
\begin{solution}
	置 \(\tau=-t\)，分离球面与轴向积分，得到
	\[
		\Phi(t)=(4\pi\tau)^{-m/2}
		\sigma_k(2k\tau)^{k/2}e^{-k/2}
		(4\pi\tau)^{(m-k)/2}
		=\sigma_k\left(\frac{k}{2\pi e}\right)^{k/2}.
	\]
	时间因子抵消，轴向高斯积分也恰抵消相应归一化因子。所得常数符合圆柱自收缩方程与单调公式的等号条件。例如 \(k=1\) 时为 \(\sqrt{2\pi/e}\)，\(k=2\) 时为 \(4/e\)。
\end{solution}

\begin{exercise}\label{g3:ex:mcf-shrinker-radius}
	设 \(X:\Sigma^m\to\mathbb R^N\) 为闭自收缩浸入，\(m\ge1\)。证明
	\[
		\int_\Sigma|X^\perp|^2\,\mu_g=2m\operatorname{Vol}_g(\Sigma),
		\qquad \max_\Sigma|X|\ge\sqrt{2m}.
	\]
	若整个像位于以原点为中心、半径 \(\sqrt{2m}\) 的闭球内，证明 \(|X|\equiv\sqrt{2m}\)，并判断它作为该球面的子流形是否极小。
\end{exercise}
\begin{solution}
	由 \(-\Delta X=\mathcal H\) 及 \(|\mathrm dX|^2=m\)，
	\[
		-\Delta|X|^2=2m+2\langle X,\mathcal H\rangle
		=2m-|X^\perp|^2.
	\]
	闭流形上积分即得恒等式，也给出最大半径下界。若 \(|X|^2\le2m\)，则 \(|X^\perp|^2\le|X|^2\le2m\)，积分等号迫使两处逐点等号，因此 \(X=X^\perp\)、\(|X|^2=2m\)。半径 \(r\) 球面的 Euclid 第二基本形式为 \(-gX/r^2\)，所以浸入的 Euclid 曲率迹与球内曲率迹关系为 \(\mathcal H_{\mathbb R^N}=\mathcal H_{S^{N-1}(r)}-mX/r^2\)。取 \(r^2=2m\) 并用自收缩方程，得到球内曲率迹为零，即它在球面中极小。
\end{solution}

\begin{exercise}\label{g3:ex:mcf-level-radial}
	在 \(\mathbb R^{m+1}\) 上取 \(u(x,t)=|x|^2+2mt-R^2\)，\(R>0\)。验证它满足水平集平均曲率方程，包括 \(x=0\) 的粘性条件。描述零水平集在 \(t<R^2/(2m)\)、\(t=R^2/(2m)\) 与之后的形状。
\end{exercise}
\begin{solution}
	当 \(x\ne0\)，有 \(\nabla u=2x\)、\(D^2u=2I\)，所以
	\((I-\nabla u\otimes\nabla u/|\nabla u|^2):D^2u=2m=u_t\)。在 \(x=0\)，\(\operatorname{tr}(2I)-\lambda_{\min}(2I)=\operatorname{tr}(2I)-\lambda_{\max}(2I)=2m\)。从上接触的测试函数与 \(u\) 有相同时间导数、空间梯度，且空间 Hessian 至少为 \(2I\)；算子上包络对 Hessian 单调，故次解条件成立。从下接触同理给超解条件。零集依次是半径 \(\sqrt{R^2-2mt}\) 的球面、单点 \(\{0\}\) 与空集。此处直接验证局部方程，没有对无界初值套用紧集外常值的存在定理。
\end{solution}

\begin{exercise}\label{g3:ex:mcf-level-relabel}
	设 \(Q\subset\mathbb R^n\times(0,T)\) 开，\(n\ge2\)，\(u:Q\to\mathbb R\) 是空间二阶、时间一阶光滑的水平集粘性解。证明 \(v=u^3\) 仍为水平集粘性解。若 \(u\) 在零水平集附近梯度非零，比较两种描述函数的零集及其上的梯度，说明梯度退化是否必然表示肥化。
\end{exercise}
\begin{solution}
	在 \(u\ne0\) 的开集，函数 \(z\mapsto z^3\) 的导数为正，故\cref{g3:prop:mcf-level-set-relabel}给出结论。在 \(u=0\) 的点，直接微分得 \(v_t=0\)、\(\mathrm Dv=0\)、\(\mathrm D^2v=0\)。从上接触的测试函数满足时间导数为零、空间梯度为零、Hessian 非负；其 \(\mathcal F^*(0,X)\) 是最大的 \(n-1\) 个特征值之和，故非负，次解条件成立。从下接触时 Hessian 非正，其 \(\mathcal F_*(0,X)\le0\)，给出超解条件。于是 \(v\) 在整个区域求解。
	两函数的零集完全相同。原来的零集由隐函数定理是光滑超曲面，故无内点，但 \(v\) 在其上处处梯度为零。因此描述函数的梯度退化并不必然表示零水平集肥化。
\end{solution}

\begin{exercise}\label{g3:ex:mcf-brakke-weight-jumps}
	设 \(F_t:M^m\to\mathbb R^N\) 为闭光滑嵌入平均曲率流，\(0\le t<T\)，\(\mu_t\) 为其面积测度。令 \(a(t)\) 为右连续、非负整数值、在每个紧时间区间只有有限次跳跃的非增函数。证明 \(\widetilde\mu_t=a(t)\mu_t\) 满足积分形式的 Brakke 不等式。它如何改变支集与重数？
\end{exercise}
\begin{solution}
	在 \(a\) 为常值正整数的区间，整变分流形的曲率仍为原来的 \(\mathcal H\)，光滑输运等式两边同乘 \(a\)。在跳跃时刻 \(c\)，非负测试函数的积分产生附加项
	\((a(c)-a(c-))\mu_c(\phi_c)\le0\)。分段积分并加上全部跳跃项，恰得 Brakke 不等式；\(a=0\) 的区间两边都为零。时间端点上的下降跳跃同样只增加非正项。正权期间支集仍是 \(F_t(M)\)，但面积重数变成 \(a(t)\)；降到零以后支集消失。故支集本身不记录正整数重数及其下降。
\end{solution}

\begin{exercise}\label{g3:ex:mcf-boundary-varifold}
	设 \(D\) 是 \(\mathbb R^2\subset\mathbb R^3\) 内的闭单位圆盘，把 \(\mathcal H^2\llcorner D\) 作为重数一变分流形的权测度。证明它虽然内点曲率为零，却不能以静止族直接满足本章无边界 Brakke 流的定义。
\end{exercise}
\begin{solution}
	对紧支环境向量场 \(X\)，平面散度定理给
	\[
		\delta V(X)=\int_D\operatorname{div}_{\mathbb R^2}X\,\mathrm d\mathcal H^2
		=\int_{\partial D}\langle X,\eta\rangle\,\mathrm d\mathcal H^1.
	\]
	右边是支于圆周的非零向量测度，而面积测度赋予圆周零质量。因此第一变分不对面积测度绝对连续，不能写成 \(-\mathcal H\mu\) 且无奇异部分。把内点曲率零代入演化不等式会漏掉边界力。本章的无边界定义因而不包含这个静止圆盘；带边界弱流需要另行规定边界项。
\end{solution}
